\documentclass[11pt,a4paper]{article}

\usepackage[T1]{fontenc}
\usepackage[utf8]{inputenc}
\usepackage[ngerman]{babel}
\usepackage{lmodern}
\usepackage{microtype}
\usepackage{geometry}
\usepackage{amsmath,amssymb,amsthm,mathtools}
\usepackage{tikz}
\usepackage{hyperref}
\usepackage{enumitem}

\geometry{margin=28mm}
\hypersetup{colorlinks=true,linkcolor=blue,citecolor=blue,urlcolor=blue}

\newcommand{\N}{\mathbb{N}}
\newcommand{\R}{\mathbb{R}}
\newcommand{\Q}{\mathbb{Q}}
\newcommand{\asympnote}{\mathrel{\asymp}}
\newcommand{\Logit}{\operatorname{logit}}
\newcommand{\diam}{\operatorname{diam}}
\newcommand{\Haus}{\operatorname{dim}_{H}}

\newtheorem{remark}{Bemerkung}[section]
\newtheorem{question}[remark]{Arbeitsfrage}
\newtheorem{conjecture}[remark]{Vermutung}
\newtheorem{proposition}[remark]{Proposition}
\newtheorem{corollary}[remark]{Korollar}

\title{Von gleichmäßigen Punkten zu projektiver Tangenzdynamik\\
\large Eine explorative Forschungsnotiz auf der Strecke $[0,1]$}
\author{Julian L.\ Voigt}
\date{15. Juli 2026}

\begin{document}
\maketitle

\begin{abstract}
Aus einer elementaren Frage über gleichmäßig angeordnete Punkte auf einer
Strecke entstehen zwei verschiedene Verfeinerungsgeometrien. Die globale
Halbierung liefert die additiv gleichmäßige dyadische Zerlegung. Eine
Tangenzkonfiguration von Kreisen liefert dagegen nach Projektion auf
$[0,1]$ eine projektiv gleichmäßige Möbius-Bahn. Diese Notiz fasst die
beiden Modelle, ihre symbolische Stern--Brocot-Dynamik, eine zugeordnete
Lückenzeta-Funktion und erste multifraktale Fragen zusammen. Sie trennt
bewusst etablierte Rechnungen von Aussagen, deren vollständiger Beweis noch
aussteht.
\end{abstract}

\paragraph{Status.}
Dies ist ein Forschungsentwurf und kein abgeschlossenes Paper. Aussagen mit
\emph{bedingt} setzen die in Abschnitt~\ref{sec:dictionary} formulierte
geometrisch-symbolische Vergleichsannahme voraus. Als \emph{Vermutung}
bezeichnete Aussagen sind durch die bisherige Exploration motiviert, aber
noch nicht bewiesen.

\section{Ausgangspunkt: Gleichmäßigkeit auf einer Strecke}

Wir betrachten die abgeschlossene Strecke $[0,1]$. Die Endpunkte $0,1$ sind
fest; zunächst ist auch der Mittelpunkt gesetzt:
\[
  P_0=\left\{0,\frac12,1\right\}.
\]
Die naive Frage lautet: Wie lassen sich Punkte hinzufügen, sodass alle
benachbarten Punkte denselben Abstand haben? Entscheidend ist, welche Metrik
das Wort ,,derselbe Abstand'' meint.

\subsection{Additive Halbierung}

Die strengste lokale Regel setzt in jedem Zug den Mittelpunkt jeder
vorhandenen Teilstrecke. Nach $k$ Zügen gilt daher
\[
 P_k=\left\{\frac{j}{2^{k+1}}:j=0,\ldots,2^{k+1}\right\},
 \qquad d_k=2^{-(k+1)}.
\]
Damit sind genau die dyadischen Brüche erreichbar. Die in Schritt $k$ neu
entstehende Generation besteht aus den Punkten mit ungeradem Zähler im
Nenner $2^{k+1}$. Insbesondere sind die Raster verschachtelt,
$P_m\subseteq P_n$ für $m\le n$, und $\bigcup_{k\ge0}P_k$ ist abzählbar
und dicht in $[0,1]$.

Die dyadische Konstruktion verteilt nicht nur Punkte, sondern auch
Information: Nach $k$ Schritten bestimmt ein Wort der Länge $k+1$ in
$\{0,1\}$ ein Teilintervall.

\section{Projektive Gleichmäßigkeit aus Tangenz}

Die zweite Verteilung entsteht aus einer symmetrischen Skelettkette von
Kreisen. Ihre Ausgangsgeometrie ist die abgeschlossene obere Halbebene des
Einheitskreises
\[
 \mathcal U=\{(x,y)\in\mathbb R^2:x^2+y^2\le1,\ y\ge0\}
\]
mit Durchmesser $[-1,1]\times\{0\}$. Ein Skelettkreis ist ein Kreis im
Inneren von $\mathcal U$, der sowohl den Durchmesser als auch den
Einheitskreis berührt. Sei $(x,0)$ sein Berührpunkt mit dem Durchmesser und
normiere diesen durch $p=(x+1)/2\in(0,1)$.

\begin{proposition}[Rekonstruktion eines Skelettkreises]
Der Skelettkreis mit Parameter $p$ besitzt Mittelpunkt und Radius
\begin{equation}\label{eq:circle-reconstruction}
 p\longmapsto \bigl((2p-1,2p(1-p)),\,2p(1-p)\bigr),
 \qquad r(p)=2p(1-p),
\end{equation}
wobei der erste Eintrag der Kreismittelpunkt und der zweite der Radius ist.
\end{proposition}

\begin{proof}
Da der Kreis den Durchmesser berührt, hat sein Mittelpunkt die Form $(x,r)$.
Die innere Berührung mit dem Einheitskreis ist äquivalent zu
\[
 x^2+r^2=(1-r)^2.
\]
Folglich ist $r=(1-x^2)/2$. Mit $x=2p-1$ folgt unmittelbar
$r=2p(1-p)$ und damit \eqref{eq:circle-reconstruction}.
\end{proof}

\begin{figure}[htbp]
 \centering
 \begin{tikzpicture}[x=4.1cm,y=4.1cm,baseline=(current bounding box.center)]
  \colorlet{outer}{black!70}
  \colorlet{circle}{blue!60!black}
  \colorlet{circlefill}{blue!12}
  \colorlet{guide}{black!35}

  \draw[outer,very thick] (-1,0) arc[start angle=180,end angle=0,radius=1];
  \draw[outer,thick] (-1.08,0) -- (1.08,0);

  \foreach \x/\r in {0/0.5,{0.70710678}/0.25,{-0.70710678}/0.25,{0.94280904}/0.05555556,{-0.94280904}/0.05555556} {
    \draw[circle,fill=circlefill] (\x,\r) circle[radius=\r];
    \draw[guide,densely dashed] (\x,0) -- (\x,\r);
    \fill[circle] (\x,0) circle[radius=0.012];
  }

  \node[font=\scriptsize,below] at (0,0) {$p_0=\tfrac12$};
  \node[font=\scriptsize,below left] at (-0.70710678,0) {$p_{-1}$};
  \node[font=\scriptsize,below right] at (0.70710678,0) {$p_1$};
  \node[font=\scriptsize,below left] at (-1,0) {$p=0$};
  \node[font=\scriptsize,below right] at (1,0) {$p=1$};

  \node[font=\small,align=center] at (0,1.16)
    {Einheitskreis $x^2+y^2=1$ und Skelettkreise};
  \node[font=\scriptsize,align=center] at (0,-0.25)
    {$p=(x+1)/2$, \quad $r(p)=2p(1-p)$};
\end{tikzpicture}

 \caption{Die ersten Skelettkreise im Einheitskreis. Ihre Berührpunkte auf
 dem Durchmesser werden durch $p=(x+1)/2$ auf $[0,1]$ normiert.}
 \label{fig:skeleton-projection}
\end{figure}

Setze
\[
 \lambda=(1+\sqrt2)^2=3+2\sqrt2.
\]
\begin{proposition}[Projektive Tangenzbedingung]
Für $0<p<q<1$ sind die Skelettkreise zu $p$ und $q$ genau dann äußerlich
tangential, wenn
\begin{equation}\label{eq:crossratio}
 \frac{q(1-p)}{p(1-q)}=\lambda.
\end{equation}
Insbesondere erfüllen aufeinanderfolgende projizierte Skelettpunkte
\begin{equation}\label{eq:mobius}
 p_{n+1}=T_\lambda(p_n)
 =\frac{\lambda p_n}{1+(\lambda-1)p_n}.
\end{equation}
\end{proposition}

\begin{proof}
Für Radien $r,s$ und Durchmesserkoordinaten $x,y$ ist äußere Tangenz
äquivalent zu $(x-y)^2=4rs$. Mit \eqref{eq:circle-reconstruction} wird das
\[
 (q-p)^2=4pq(1-p)(1-q).
\]
Setze $t=p/(1-p)$ und $u=q/(1-q)$. Nach Kürzen bleibt
$(u-t)^2=4ut$. Für $u>t$ liefert $z=\sqrt{u/t}>1$ die Gleichung
$z^2-1=2z$, also $z=1+\sqrt2$ und damit $u/t=\lambda$. Das Umstellen von
\eqref{eq:crossratio} ergibt \eqref{eq:mobius}.
\end{proof}

In der Logit-Koordinate
\[
 \eta(p)=\Logit(p)=\log\frac{p}{1-p}
\]
wird \eqref{eq:mobius} zur Translation
\[
 \eta(p_{n+1})-\eta(p_n)=\log\lambda.
\]
Somit stehen die Punkte nicht im euklidischen, wohl aber im projektiven Sinn
in konstantem Abstand.

\begin{remark}
Die zentrale Gegenüberstellung lautet
\[
 \text{globale Halbierung}\quad\longleftrightarrow\quad
 \text{additive Gleichmäßigkeit},
\]
\[
 \text{Tangenzprojektion}\quad\longleftrightarrow\quad
 \text{projektive Gleichmäßigkeit}.
\]
Sie ist der eigentliche Perspektivwechsel der Ausgangsfrage.
\end{remark}

\section{Stern--Brocot-Kodierung}\label{sec:dictionary}

Zunächst betrachten wir das arithmetische Referenzmodell. Die
Randverfeinerung wird durch
\[
 L=\begin{pmatrix}1&0\\1&1\end{pmatrix},
 \qquad
 R=\begin{pmatrix}1&1\\0&1\end{pmatrix}
\]
organisiert. Der Wurzelknoten ist das Intervall $[0/1,1/1]$. Hat ein Knoten
benachbarte Endpunkte $a/u<c/v$ mit $cu-av=1$, dann setzt die Verfeinerung
den Medianten $(a+c)/(u+v)$ ein. Die Nennerpaare seiner linken und rechten
Kindintervalle sind jeweils
\[
 (u,v)\longmapsto (u,u+v)=L\binom{u}{v},
 \qquad
 (u,v)\longmapsto (u+v,v)=R\binom{u}{v}.
\]

\begin{proposition}[Exaktes Stern--Brocot-Wörterbuch]
Die endlichen Wörter $w\in\{L,R\}^*$ stehen bijektiv zu den positiven
teilerfremden Paaren $(u_w,v_w)$. Das zugehörige Referenzintervall $J_w$
hat die exakte Länge
\begin{equation}\label{eq:reference-length}
 |J_w|=\frac{1}{u_wv_w}.
\end{equation}
\end{proposition}

\begin{proof}
Die Subtraktionsversion des euklidischen Algorithmus reduziert jedes
positive teilerfremde Paar eindeutig auf $(1,1)$; rückwärts entstehen genau
die beiden angegebenen Operationen $L,R$. Für benachbarte Endpunkte gilt
\[
 \frac{c}{v}-\frac{a}{u}=\frac{cu-av}{uv}=\frac1{uv},
\]
was \eqref{eq:reference-length} beweist.
\end{proof}

Ein unendliches Wort beschreibt einen Grenzpunkt. Wie beim binären
Stellenwertsystem sind die endlich erzeugten Randpunkte die Ausnahme mit
zwei symbolischen Adressen. Für die geometrische Lücke der Tangenzfigur
bleibt folgende Brücke zu beweisen:

\begin{quote}
\textbf{Vergleichsannahme.} Zu jedem Wort $w$ gehört eine geometrische
Lücke $I_w$ mit Länge $\ell_w$. Unter der Stern--Brocot-Bijektion
$w\leftrightarrow(u_w,v_w)$ zu positiven teilerfremden Paaren existieren
Konstanten $0<c\le C<\infty$, unabhängig von $w$, so dass
\begin{equation}\label{eq:comparison}
 \frac{c}{u_wv_w}\le \ell_w\le\frac{C}{u_wv_w}.
\end{equation}
\end{quote}

Die besondere Wiederholung der ersten Lücke wird durch das Wort $LR^2L$
ausgedrückt. Sein Matrixprodukt besitzt den dominanten Eigenwert
$3+2\sqrt2$; daher tritt dieselbe Zahl sowohl in der Skelettbahn als auch
in einer Pell-artigen Rekursion
\[
 u_{k+1}=6u_k-u_{k-1}
\]
auf. Hier muss noch präzisiert werden, welche konkrete Folge $u_k$ gemeint
ist und wie sie aus der Geometrie normiert wird.

\section{Konkrete Randlücken der Halbkreisgeometrie}

Wir wenden das Wörterbuch nun auf eine geometrisch bestimmte Teilfamilie der
Packung an. Für $n\ge0$ sei $G_n$ die Lücke, die von der Halbierungsgeraden,
$C_n$ und $C_{n+1}$ begrenzt wird. Wir verfolgen nur diejenigen rekursiv
entstehenden Teillücken, die weiterhin an der Halbierungsgeraden liegen.
Sie bilden den \emph{Randbaum} von $G_n$. Die obere, nicht an der Geraden
liegende Teilpackung wird hier bewusst nicht mitgezählt.

Seien $A,B$ zwei benachbarte Kreise des Randbaums mit Radien $a^{-2}$ und
$b^{-2}$. Ihre Berührpunkte mit der Halbierungsgeraden begrenzen ein Intervall
der Länge
\begin{equation}\label{eq:gap-length-exact}
 \ell(A,B)=\frac{2}{ab}.
\end{equation}
Der neue Kreis, der an $A$, $B$ und die Halbierungsgerade tangential ist,
hat Radius $(a+b)^{-2}$.

\begin{proposition}[Exakte geometrische Randrekursion]
Die beiden Kindlücken eines Randintervalls mit Parameterpaar $(a,b)$ besitzen
die Parameterpaare
\[
 (a,a+b)=L\binom{a}{b},
 \qquad
 (a+b,b)=R\binom{a}{b}.
\]
Insbesondere gilt \eqref{eq:gap-length-exact} auf jedem Knoten des
Randbaums exakt.
\end{proposition}

\begin{proof}
Zwei an dieselbe Gerade tangentiale Kreise mit Radien $a^{-2}$ und $b^{-2}$
sind genau dann äußerlich tangential, wenn der Abstand ihrer Berührpunkte auf
der Geraden $2/(ab)$ beträgt. Für einen Kreis mit Radius $(a+b)^{-2}$ sind
die beiden Teilabstände
\[
 \frac{2}{a(a+b)}
 \quad\text{und}\quad
 \frac{2}{b(a+b)}.
\]
Ihre Summe ist $2/(ab)$; der neue Kreis passt also exakt zwischen die beiden
alten. Seine beiden Nachbarpaare sind folglich $(a,a+b)$ und $(a+b,b)$.
\end{proof}

Für die Skelettkette aus dem Paper gilt mit $q=1+\sqrt2$ und
$\alpha=\log q$
\[
 r_n=\frac{1}{2\cosh^2(n\alpha)},
 \qquad
 a_n=r_n^{-1/2},\qquad b_n=r_{n+1}^{-1/2}.
\]
Der Quotient der Wurzelparameter erfüllt gleichmäßig in $n$
\begin{equation}\label{eq:root-ratio}
 \sqrt2\le \frac{b_n}{a_n}
 =\frac{\cosh((n+1)\alpha)}{\cosh(n\alpha)}<q.
\end{equation}

Sei $(u_w,v_w)=M_w(1,1)$ das Stern--Brocot-Paar des Wortes $w$. Bezeichne
$\ell_{n,w}$ die Länge der geometrischen Randlücke in $G_n$ mit Adresse $w$.

\begin{proposition}[Uniformer Vergleich aus der Geometrie]
Für alle $n\ge0$ und $w\in\{L,R\}^*$ gilt
\begin{equation}\label{eq:geometric-comparison}
 \frac{2r_n}{q^2u_wv_w}
 \le \ell_{n,w}\le
 \frac{2r_n}{u_wv_w}.
\end{equation}
\end{proposition}

\begin{proof}
Schreibe $M_w=\left(\begin{smallmatrix}A&B\\C&D\end{smallmatrix}\right)$
mit nichtnegativen Einträgen. Der Parameter des Knotens ist
\[
 M_w\binom{a_n}{b_n}
 =a_n\binom{A+B(b_n/a_n)}{C+D(b_n/a_n)}.
\]
Nach \eqref{eq:root-ratio} liegt jede der beiden Komponenten zwischen der
entsprechenden Komponente von $M_w(1,1)$ und deren $q$-Fachem. Die exakte
Längenformel \eqref{eq:gap-length-exact} und $a_n^{-2}=r_n$ ergeben
\eqref{eq:geometric-comparison}.
\end{proof}

\section{Zeta-Funktion und Zählung der konkreten Randlücken}

Wir zählen zunächst nur die rechte Randfamilie. Ihre Zeta-Funktion ist
\[
 Z_\partial^+(s)=
 \sum_{n\ge0}\sum_{w\in\{L,R\}^*}\ell_{n,w}^s.
\]
Die gespiegelte linke Familie liefert lediglich den Faktor $2$ und hat
denselben kritischen Exponenten. Das arithmetische Referenzmodell ist
\[
 Z_{\mathrm{ref}}(s)
 =\sum_{\substack{u,v\ge1\\\gcd(u,v)=1}}\frac{1}{(uv)^s}
 =\frac{\zeta(s)^2}{\zeta(2s)},\qquad s>1.
\]
Die letzte Gleichheit folgt aus der Möbius-Inversion oder aus der
Eulerproduktzerlegung. Es besitzt bei $s=1$ einen Doppelpol:
\[
 Z_{\mathrm{ref}}(s)\sim\frac{6}{\pi^2}(s-1)^{-2}.
\]
Der Faktor $2^s$ aus einer alternativen Konvention für orientierte oder
beidseitige Lücken würde den Koeffizienten bei $s=1$ zu $12/\pi^2$
ändern. Diese Konvention muss vor einer endgültigen Aussage fixiert
werden.

\begin{proposition}[Kritischer Exponent der konkreten Randfamilie]
Die Reihe $Z_\partial^+(s)$ konvergiert genau für $s>1$. Insbesondere ist
der kritische Exponent der Randfamilie
\[
 s_c=1.
\]
\end{proposition}

\begin{proof}
Aus \eqref{eq:geometric-comparison} folgt für $s>0$
\[
 \frac{2^s}{q^{2s}}\Bigl(\sum_{n\ge0}r_n^s\Bigr)Z_{\mathrm{ref}}(s)
 \le Z_\partial^+(s)\le
 2^s\Bigl(\sum_{n\ge0}r_n^s\Bigr)Z_{\mathrm{ref}}(s).
\]
Wegen $r_n\asymp q^{-2n}$ konvergiert die erste Summe für jedes $s>0$.
Die Behauptung folgt daher aus der Eulerproduktformel für
$Z_{\mathrm{ref}}$.
\end{proof}

Für die Zählung setzen wir
\[
 N_\partial^+(\varepsilon)=
 \#\{(n,w):\ell_{n,w}\ge\varepsilon\}.
\]

\begin{proposition}[Lückenzählung der konkreten Randfamilie]
Für $\varepsilon\downarrow0$ gilt
\[
 N_\partial^+(\varepsilon)=\Theta\!\left(
 \varepsilon^{-1}\log\frac1\varepsilon\right),
\]
und dieselbe Ordnung gilt nach Hinzunahme der gespiegelten linken Familie.
\end{proposition}

\begin{proof}
Sei
\[
 F(X)=\#\{(u,v)\in\N^2:\gcd(u,v)=1,\ uv\le X\}.
\]
Möbius-Inversion in der klassischen Divisorsumme liefert
$F(X)=\Theta(X\log(2+X))$. Aus \eqref{eq:geometric-comparison} folgt
\[
 \sum_{n\ge0}F\!\left(\frac{2r_n}{q^2\varepsilon}\right)
 \le N_\partial^+(\varepsilon)\le
 \sum_{n\ge0}F\!\left(\frac{2r_n}{\varepsilon}\right).
\]
Die untere Schranke folgt bereits aus dem Summanden $n=0$. Für die obere
Schranke tragen nur $n$ mit $2r_n\ge\varepsilon$ bei. Mit der exponentiellen
Abnahme von $r_n$ und $\sum_n r_n<\infty$ ist die rechte Seite
$O(\varepsilon^{-1}\log(1/\varepsilon))$.
\end{proof}

\begin{remark}
Damit sind Exponent und Zählordnung für eine kanonische, geometrisch
definierte Teilfamilie der vollen Halbkreispackung bewiesen.
\end{remark}

\section{Die Randfamilie am Außenkreis}

Die zweite Randfamilie besitzt dieselbe Struktur. Sichtbar wird sie nach
einer Möbius-Transformation. Identifiziere die Ebene mit $\mathbb C$ und
betrachte die Cayley-Transformation
\[
 \Phi(z)=i\frac{1+z}{1-z},
 \qquad
 \Phi^{-1}(t)=\frac{t-i}{t+i}.
\]
Sie bildet den Einheitskreis ohne den Punkt $1$ auf die reelle Gerade ab.
Sei $\xi_n$ der Berührpunkt von $C_n$ mit dem Einheitskreis. Aus den
geschlossenen Skelettformeln folgt
\begin{equation}\label{eq:cayley-skeleton}
 \Phi(\xi_n)=-q^{2n},
 \qquad
 \operatorname{rad}(\Phi(C_n))=q^{2n}.
\end{equation}
Tatsächlich ist $\xi_n=(x_n+ir_n)/(1-r_n)$, und die Standardformel für den
Bildradius eines Kreises unter $\Phi$ liefert
\[
 \operatorname{rad}(\Phi(C_n))
 =\frac{2r_n}{|1-(x_n+ir_n)|^2-r_n^2}
 =\frac{2r_n}{(1-x_n)^2}=q^{2n}.
\]

Damit werden die Kreise, die in einer oberen Lücke gleichzeitig den
Außenkreis berühren, zu Kreisen über der reellen Achse. Die Randparameter
des Wurzelintervalls zwischen $C_n$ und $C_{n+1}$ sind im Bild
\[
 \widetilde a_n=q^{-n},
 \qquad
 \widetilde b_n=q^{-(n+1)}.
\]
Die exakte Rekursion aus Proposition~4.1 gilt dort unverändert.

Sei $\widetilde\ell_{n,w}$ die Länge des Bildintervalls auf der reellen
Achse und $\ell^{\circ}_{n,w}$ die Bogenlänge seines Urbilds auf dem
Einheitskreis. Dann gilt für jedes Wort $w$
\begin{equation}\label{eq:cayley-real-comparison}
 \frac{2q^{2n}}{u_wv_w}
 \le \widetilde\ell_{n,w}\le
 \frac{2q^{2n+2}}{u_wv_w}.
\end{equation}
Denn für $M_w=\left(\begin{smallmatrix}A&B\\C&D\end{smallmatrix}\right)$
liegen $A+B/q$ und $C+D/q$ jeweils zwischen der entsprechenden
Komponente von $M_w(1,1)$ dividiert durch $q$ und dieser Komponente selbst.
Die exakte Geradenformel \eqref{eq:gap-length-exact} liefert dann
\eqref{eq:cayley-real-comparison}.
Auf dem ganzen Wurzelintervall liegen die reellen Koordinaten zwischen
$-q^{2n+2}$ und $-q^{2n}$. Da
\[
 \left|\frac{d}{dt}\Phi^{-1}(t)\right|=\frac{2}{1+t^2},
\]
folgt
\begin{equation}\label{eq:outer-arc-comparison}
 \frac{4q^{2n}}{(1+q^{4n+4})u_wv_w}
 \le \ell^{\circ}_{n,w}\le
 \frac{4q^{2n+2}}{(1+q^{4n})u_wv_w}.
\end{equation}

\begin{proposition}[Außenkreisfamilie]
Für die rechte Außenkreisfamilie definieren wir mit ihren Bogenlängen
\[
 Z_\circ^+(s)=\sum_{n\ge0}\sum_{w\in\{L,R\}^*}
 (\ell^{\circ}_{n,w})^s
\]
gilt erneut $s_c=1$. Ferner ist
\[
 N_\circ^+(\varepsilon)=
 \Theta\!\left(\varepsilon^{-1}\log\frac1\varepsilon\right).
\]
\end{proposition}

\begin{proof}
Aus \eqref{eq:outer-arc-comparison} folgt mit Konstanten, die nur von $q$
abhängen,
\[
 \ell^{\circ}_{n,w}\asymp\frac{q^{-2n}}{u_wv_w}.
\]
Die Argumente aus Propositionen~5.1 und~5.2 gelten daher mit der
geometrischen Reihe $\sum_n q^{-2ns}$ unverändert.
\end{proof}

\begin{remark}
Beide Randfamilien der Halbkreispackung besitzen damit dieselbe kritische
Skalierung. Offen bleibt jetzt nur noch die Zeta-Analyse der nicht an einen
Rand gebundenen, zweidimensionalen Apollonischen Teilpackungen.
\end{remark}

\section{Die erste Binnenebene}

Die Randbäume erzeugen bereits eine kanonische Teilfamilie von Kreisen, die
nicht mehr an einem Rand liegen. Sei $G_{n,w}$ eine untere Randlücke mit
Parameterpaar $(a,b)$. Nach dem Einsetzen des geradentangentialen Kreises mit
Krümmung $(a+b)^2$ entsteht eine Dreieckslücke. Ihr erster Binnenkreis ist
der zweite Descartes-Kreis zu dieser Konfiguration. Seine Krümmung ist
\[
 \kappa(a,b)=2\bigl(a^2+b^2+(a+b)^2\bigr)
 =4(a^2+ab+b^2),
\]
also sein Radius
\begin{equation}\label{eq:seed-radius-exact}
 \rho(a,b)=\frac{1}{4(a^2+ab+b^2)}.
\end{equation}
Wir nennen die so entstehenden Kreise $E_{n,w}$ die \emph{erste
Binnenebene}.

\begin{proposition}[Uniformer Vergleich der ersten Binnenebene]
Sei $(u_w,v_w)=M_w(1,1)$ und sei $\rho_{n,w}$ der Radius von $E_{n,w}$.
Dann gilt für alle $n\ge0$ und $w\in\{L,R\}^*$
\begin{equation}\label{eq:seed-comparison}
 \frac{r_n}{4q^2(u_w+v_w)^2}
 \le \rho_{n,w}\le
 \frac{r_n}{3(u_w+v_w)^2}.
\end{equation}
\end{proposition}

\begin{proof}
Schreibe $(a,b)=M_w(a_n,b_n)$ und $\theta=b_n/a_n$. Aus
\eqref{eq:root-ratio} folgt
\[
 a_nu_w\le a\le q a_nu_w,
 \qquad
 a_nv_w\le b\le q a_nv_w.
\]
Weiter gilt
\[
 \frac34(u_w+v_w)^2\le u_w^2+u_wv_w+v_w^2
 \le (u_w+v_w)^2.
\]
Setzt man diese Abschätzungen in \eqref{eq:seed-radius-exact} ein und
verwendet $a_n^{-2}=r_n$, so erhält man \eqref{eq:seed-comparison}.
\end{proof}

Die dazugehörige Referenzzeta ist nicht mehr die Produktreihe aus dem
Randbaum, sondern
\[
 Z_{\mathrm{seed,ref}}(s)
 =\sum_{\substack{u,v\ge1\\\gcd(u,v)=1}}\frac{1}{(u+v)^{2s}}
 =\sum_{m\ge2}\frac{\varphi(m)}{m^{2s}}
 =\frac{\zeta(2s-1)}{\zeta(2s)}
 \qquad(s>1).
\]
Sie hat bei $s=1$ einen einfachen Pol. Die andere Arithmetik zeigt sich
auch in der Zählung.

\begin{proposition}[Zeta und Zählung der ersten Binnenebene]
Für
\[
 Z_{\mathrm{seed}}^+(s)=
 \sum_{n\ge0}\sum_{w\in\{L,R\}^*}\rho_{n,w}^s
\]
gilt der kritische Exponent $s_c=1$ und, für $s\downarrow1$,
\[
 Z_{\mathrm{seed}}^+(s)=\Theta\!\left(\frac{1}{s-1}\right).
\]
Ferner erfüllt die Radiuszählfunktion
\[
 N_{\mathrm{seed}}^+(\varepsilon)=
 \#\{(n,w):\rho_{n,w}\ge\varepsilon\}
\]
die Ordnung
\[
 N_{\mathrm{seed}}^+(\varepsilon)=\Theta(\varepsilon^{-1})
 \qquad(\varepsilon\downarrow0).
\]
\end{proposition}

\begin{proof}
Die Zeta-Aussage folgt aus \eqref{eq:seed-comparison}, der Konvergenz von
$\sum_n r_n^s$ für $s>0$ und der angegebenen Dirichlet-Reihe für
$Z_{\mathrm{seed,ref}}$. Für die Zählung setze
\[
 H(X)=\#\{(u,v)\in\N^2:\gcd(u,v)=1,\ u+v\le X\}.
\]
Wegen $H(X)=\sum_{m\le X}\varphi(m)=\Theta(X^2)$ liefert
\eqref{eq:seed-comparison} nach Summation über $n$ die obere Schranke
$O(\varepsilon^{-1}\sum_n r_n)=O(\varepsilon^{-1})$. Die untere Schranke
folgt bereits aus $n=0$ und $H(c\varepsilon^{-1/2})=\Theta(\varepsilon^{-1})$.
\end{proof}

\begin{remark}
Die erste Binnenebene hat also noch denselben kritischen Exponenten wie die
Randfamilie, aber keinen logarithmischen Zählfaktor. Dieser Unterschied
trennt die Produktarithmetik $uv$ der Randlücken von der Summenarithmetik
$(u+v)^2$ der ersten freien Descartes-Kreise. Eine Aussage über alle tieferen
Binnenkreise ist damit ausdrücklich noch nicht getroffen.
\end{remark}

\subsection{Zweiseitige Leitkoeffizienten}

Setze
\[
 R(s)=\sum_{n\ge0}r_n^s,
 \qquad R_1=R(1).
\]
Wegen $r_n\asymp q^{-2n}$ ist $R(s)$ für jedes $s>0$ endlich und
$R(s)\to R_1$ für $s\downarrow1$. Die folgenden Aussagen bestimmen
noch nicht die exakten geometrischen Leitkoeffizienten. Sie machen aber
die aus dem uniformen Vergleich bereits erzwungene Koeffizientenhülle
explizit.

\begin{corollary}[Abel-Hülle der Randzeta]
Für die rechte Randfamilie gilt
\[
 \frac{12R_1}{\pi^2q^2}
 \le
 \liminf_{s\downarrow1}(s-1)^2Z_\partial^+(s)
 \le
 \limsup_{s\downarrow1}(s-1)^2Z_\partial^+(s)
 \le
 \frac{12R_1}{\pi^2}.
\]
\end{corollary}

\begin{proof}
Der Vergleich \eqref{eq:geometric-comparison} liefert für $s>1$
\[
 \frac{2^s}{q^{2s}}R(s)Z_{\mathrm{ref}}(s)
 \le Z_\partial^+(s)
 \le 2^sR(s)Z_{\mathrm{ref}}(s).
\]
Da
\[
 Z_{\mathrm{ref}}(s)
 \sim\frac6{\pi^2}(s-1)^{-2}
 \qquad(s\downarrow1),
\]
folgt die Behauptung nach Multiplikation mit $(s-1)^2$ und
Grenzübergang.
\end{proof}

\begin{corollary}[Leitkoeffizientenhülle der Randzählung]
Für die rechte Randfamilie gilt
\[
 \frac{12R_1}{\pi^2q^2}
 \le
 \liminf_{\varepsilon\downarrow0}
 \frac{\varepsilon N_\partial^+(\varepsilon)}
 {\log(1/\varepsilon)}
 \le
 \limsup_{\varepsilon\downarrow0}
 \frac{\varepsilon N_\partial^+(\varepsilon)}
 {\log(1/\varepsilon)}
 \le
 \frac{12R_1}{\pi^2}.
\]
\end{corollary}

\begin{proof}
Für
\[
 F(X)=\#\{(u,v)\in\mathbb N^2:\gcd(u,v)=1,\ uv\le X\}
\]
gilt die klassische primitive Divisorasymptotik
\[
 F(X)\sim\frac6{\pi^2}X\log X.
\]
Der Vergleich \eqref{eq:geometric-comparison} ergibt
\[
 \sum_{n\ge0}F\!\left(\frac{2r_n}{q^2\varepsilon}\right)
 \le N_\partial^+(\varepsilon)
 \le
 \sum_{n\ge0}F\!\left(\frac{2r_n}{\varepsilon}\right).
\]
Für jede feste endliche Teilsumme folgt die Behauptung aus der
Asymptotik von $F$. Für die Restsummen gilt
$F(X)\ll X\log(2+X)$. Nach Multiplikation mit
$\varepsilon/\log(1/\varepsilon)$ werden sie daher durch ein
konstantes Vielfaches von $\sum_{n>N}r_n$ kontrolliert. Zuerst
$\varepsilon\downarrow0$ und dann $N\to\infty$ liefert die
Behauptung.
\end{proof}

\begin{corollary}[Abel-Hülle der ersten Binnenebene]
Für die erste Binnenebene gilt
\[
 \frac{3R_1}{4q^2\pi^2}
 \le
 \liminf_{s\downarrow1}(s-1)Z_{\mathrm{seed}}^+(s)
 \le
 \limsup_{s\downarrow1}(s-1)Z_{\mathrm{seed}}^+(s)
 \le
 \frac{R_1}{\pi^2}.
\]
\end{corollary}

\begin{proof}
Aus \eqref{eq:seed-comparison} folgt
\[
 \frac{R(s)}{(4q^2)^s}Z_{\mathrm{seed,ref}}(s)
 \le Z_{\mathrm{seed}}^+(s)
 \le \frac{R(s)}{3^s}Z_{\mathrm{seed,ref}}(s).
\]
Wegen
\[
 Z_{\mathrm{seed,ref}}(s)
 =\frac{\zeta(2s-1)}{\zeta(2s)}
 \sim\frac3{\pi^2}(s-1)^{-1}
 \qquad(s\downarrow1)
\]
folgt die Behauptung.
\end{proof}

\begin{corollary}[Leitkoeffizientenhülle der Seed-Zählung]
Für die erste Binnenebene gilt
\[
 \frac{3R_1}{4q^2\pi^2}
 \le
 \liminf_{\varepsilon\downarrow0}
 \varepsilon N_{\mathrm{seed}}^+(\varepsilon)
 \le
 \limsup_{\varepsilon\downarrow0}
 \varepsilon N_{\mathrm{seed}}^+(\varepsilon)
 \le
 \frac{R_1}{\pi^2}.
\]
\end{corollary}

\begin{proof}
Mit
\[
 H(X)=\#\{(u,v)\in\mathbb N^2:\gcd(u,v)=1,\ u+v\le X\}
 =\sum_{m\le X}\varphi(m)
\]
gilt
\[
 H(X)\sim\frac3{\pi^2}X^2.
\]
Aus \eqref{eq:seed-comparison} folgt
\[
 \sum_{n\ge0}
 H\!\left(\sqrt{\frac{r_n}{4q^2\varepsilon}}\right)
 \le N_{\mathrm{seed}}^+(\varepsilon)
 \le
 \sum_{n\ge0}
 H\!\left(\sqrt{\frac{r_n}{3\varepsilon}}\right).
\]
Die feste Kopf--Schwanz-Zerlegung wie im Randfall und die Summierbarkeit
von $\sum_nr_n$ liefern die Behauptung.
\end{proof}

\section{Zählung der vollständigen Packung}

Die vollständige Halbkreispackung ist keine einzelne Apollonische Bahn. Sie
ist die disjunkte Vereinigung der Kreise, die in den offenen Lücken zwischen
aufeinanderfolgenden Skelettkreisen entstehen. Setze
\[
 \mathcal S^+=\{C_n:n\ge0\},
 \qquad
 \mathcal S^-=\{C_{-n}:n\ge1\}.
\]
Damit gehört der mittlere Skelettkreis $C_0$ nur zu $\mathcal S^+$.
Sei
$\mathcal A_n^{\mathrm{G}}$ beziehungsweise
$\mathcal A_n^{\mathrm{K}}$ die Menge der neu geborenen Kreise in der
Geraden- beziehungsweise Außenkreislücke zwischen $C_n$ und $C_{n+1}$.
Die Randkreise der Lücke selbst gehören nicht zu $\mathcal A_n^{\bullet}$.

\begin{proposition}[Kanonische Zählkonvention]
Mit $\widetilde{\mathcal A}_n^{\bullet}$ als linker Spiegelung von
$\mathcal A_n^{\bullet}$ ist die Menge aller Kreise der Packung die
disjunkte Vereinigung
\[
 \mathcal P=\mathcal S^+\,\dot\cup\,\mathcal S^-
 \,\dot\cup\!
 \bigsqcup_{n\ge0}\bigl(
 \mathcal A_n^{\mathrm{G}}\dot\cup\mathcal A_n^{\mathrm{K}}\bigr)
 \,\dot\cup\!
 \bigsqcup_{n\ge0}\bigl(
 \widetilde{\mathcal A}_n^{\mathrm{G}}
 \dot\cup\widetilde{\mathcal A}_n^{\mathrm{K}}\bigr).
\]
Insbesondere zählt die Radiuszeta
\[
 Z_{\mathcal P}(s)=\sum_{C\in\mathcal P}\operatorname{rad}(C)^s
\]
jeden Kreis genau einmal.
\end{proposition}

\begin{proof}
Die beiden Skelettfamilien sind wegen der getrennten Behandlung von $C_0$
disjunkt. Die offenen Geraden- und Außenkreislücken haben paarweise
disjunkte Innere; zwei benachbarte Lücken teilen höchstens Skelettkreise,
die aus den Mengen $\mathcal A_n^{\bullet}$ ausgeschlossen sind. Innerhalb
einer festen Lücke ist jeder neue Kreis der eindeutige Descartes-Kreis einer
eindeutigen Elternlücke. Damit besitzt er eine eindeutige Geburtsadresse.
\end{proof}

\begin{figure}[htbp]
 \centering
 \begin{tikzpicture}[scale=0.9, line cap=round, line join=round]
  \draw[thick] (-3,0) arc (180:0:3);
  \draw[thick] (-3,0) -- (3,0);

  \draw[thick, fill=blue!8] (0,1.5) circle (1.5);
  \draw[thick, fill=blue!16] (2.121,0.75) circle (0.75);

  \node[below] at (0,0) {$C_n$};
  \node[below] at (2.121,0) {$C_{n+1}$};
  \node[below left] at (-3,0) {$-1$};
  \node[below right] at (3,0) {$1$};

  \draw[->] (0.95,-0.95) -- (0.95,-0.08);
  \node[align=center] at (0.95,-1.32) {Geradenl\"ucke\\$\mathcal A_n^{\mathrm G}$};

  \draw[->] (2.20,2.38) -- (1.78,1.60);
  \node[align=center] at (2.35,2.78) {Au\ss enkl\"ucke\\$\mathcal A_n^{\mathrm K}$};

  \node at (-1.90,2.05) {Einheitskreis};
\end{tikzpicture}

 \caption{Lokale Zerlegung zwischen zwei aufeinanderfolgenden
 Skelettkreisen. Die Geraden- und die Außenlücke tragen jeweils eine
 eigene Apollonische Teilpackung; die Skelettkreise selbst werden nur in
 $\mathcal S$ gezählt.}
 \label{fig:packing-decomposition}
\end{figure}

Für eine einzelne normalisierte Apollonische Dreieckspackung ist die
Situation bekannt. Ihre Radiuszählung hat den Exponenten
\[
 \delta=1.305686728\ldots,
\]
die Hausdorff-Dimension des klassischen Apollonischen Gaskets. Genauer gilt
für eine solche lokale Packung $\mathcal A_\theta$, deren Formparameter
$\theta$ in einem kompakten positiven Intervall liegt,
\[
 N_\theta(\varepsilon)
 \sim c(\theta)\varepsilon^{-\delta},
 \qquad
 \lim_{s\downarrow\delta}(s-\delta)Z_\theta(s)
 =\delta c(\theta).
\]
Die Zählaussage ist der Kreis-Zählsatz von
\href{https://arxiv.org/abs/0811.2236}{Kontorovich--Oh}; die präzise
Dimension wurde jüngst von
\href{https://arxiv.org/abs/2406.04922}{Vytnova--Wormell} rigoros
berechnet.

In unserer unteren Familie ist die Normalisierung besonders transparent:
Nach Skalierung mit $r_n^{-1}$ besitzt die Wurzellücke den Parameter
\[
 \theta_n=\frac{b_n}{a_n}\in[\sqrt2,q].
\]

\begin{proposition}[Uniforme Möbius-Normierung]
Sei $\mathcal G_\theta$ die untere Apollonische Packung, deren beiden
Randkreise Radien $1$ und $\theta^{-2}$ besitzen. Für
$\theta\in[\sqrt2,q]$ bildet
\[
 M_\theta(z)=\frac{z}{1+\frac{\theta-2}{2}z}
\]
die Packung $\mathcal G_2$ auf $\mathcal G_\theta$ ab. Es gibt Konstanten
$0<m\le M<\infty$, unabhängig von $\theta$ und vom Kreis $C\in\mathcal G_2$,
so dass
\[
 m\operatorname{rad}(C)
 \le\operatorname{rad}(M_\theta(C))
 \le M\operatorname{rad}(C).
\]
\end{proposition}

\begin{proof}
Die Berührpunkte der beiden Randkreise von $\mathcal G_2$ sind $0$ und $1$.
Es gilt
\[
 M_\theta(0)=0,
 \qquad M_\theta(1)=\frac2\theta,
 \qquad M_\theta'(0)=1,
 \qquad M_\theta'(1)=\frac4{\theta^2}.
\]
Damit werden die beiden Horokreise mit Radien $1$ und $1/4$ exakt auf die
Horokreise mit Radien $1$ und $\theta^{-2}$ abgebildet; die gesamte
Descartes-Rekursion wird mitgenommen. Die Pole der Familie $M_\theta$ liegen
gleichmäßig außerhalb einer festen kompakten Umgebung aller
$\overline{\mathcal G_2}$. Auf dieser Umgebung sind $M_\theta$ und ihre
Inversen daher gleichmäßig Lipschitz. Der Vergleich der Kreisradien folgt
aus dem Vergleich der Durchmesser.
\end{proof}

\begin{figure}[htbp]
 \centering
 \begin{tikzpicture}[x=1cm,y=1cm, line cap=round, line join=round]
  \begin{scope}
    \draw[thick] (0,0) -- (3.4,0);
    \draw[thick, fill=blue!10] (0.65,0.65) circle (0.65);
    \draw[thick, fill=blue!18] (1.95,0.1625) circle (0.1625);
    \node[below] at (0.65,0) {$0$};
    \node[below] at (1.95,0) {$1$};
    \node[above] at (0.65,1.38) {$1$};
    \node[right] at (2.14,0.18) {$1/4$};
    \node[below] at (1.55,-0.55) {$\mathcal G_2$};
  \end{scope}

  \draw[->, very thick] (4.0,0.75) -- (5.95,0.75);
  \node[align=center] at (4.98,1.32) {$M_\theta(z)=\dfrac{z}{1+\frac{\theta-2}{2}z}$};

  \begin{scope}[xshift=6.6cm]
    \draw[thick] (0,0) -- (3.4,0);
    \draw[thick, fill=blue!10] (0.65,0.65) circle (0.65);
    \draw[thick, fill=blue!18] (2.00,0.25) circle (0.25);
    \node[below] at (0.65,0) {$0$};
    \node[below] at (2.00,0) {$2/\theta$};
    \node[above] at (0.65,1.38) {$1$};
    \node[right] at (2.28,0.28) {$\theta^{-2}$};
    \node[below] at (1.55,-0.55) {$\mathcal G_\theta$};
  \end{scope}
\end{tikzpicture}

 \caption{Schematische uniforme Normierung der unteren Wurzellücken.
 Die rationale Abbildung $M_\theta$ nimmt die Referenzpackung
 $\mathcal G_2$ exakt auf $\mathcal G_\theta$ mit und kontrolliert die
 Kreisradien auf dem kompakten Parameterintervall einheitlich.}
 \label{fig:moebius-normalization}
\end{figure}

Als direkte Folge gelten für die lokale Zeta und Zählung gleichmäßige
Vergleiche mit der festen Packung $\mathcal G_2$. Insbesondere ist der
kritische Exponent jeder Wurzellücke $\delta$, und die Konstanten in den
groben Abschätzungen $N_\theta(\varepsilon)\asymp
\varepsilon^{-\delta}$ können auf $[\sqrt2,q]$ einheitlich gewählt werden.

Für die Außenkreisfamilie führt die Cayley-Transformation auf eine feste
normalisierte Horokreispackung mit Parameter $q$. Nach der Skalierung
$t=-q^{2n}u$ ist ihre Rückabbildung
\[
 F_n(u)=\Phi^{-1}(-q^{2n}u)
\]
auf dem festen kompakten Wurzeldreieck gleichmäßig bi-Lipschitz bis auf den
Faktor $q^{-2n}$, denn
\[
 |F_n'(u)|=\frac{2q^{2n}}{|q^{2n}u-i|^2}\asymp q^{-2n}.
\]
Auch dort sind die lokalen Radiuszählungen also gleichmäßig mit einer festen
Apollonischen Packung vergleichbar.

\begin{proposition}[Globaler Exponent und globale Zählordnung]
Die Radiuszeta der vollständigen Halbkreispackung besitzt den kritischen
Exponenten
\[
 s_c(\mathcal P)=\delta.
\]
Ferner gilt
\[
 N_{\mathcal P}(\varepsilon)=\Theta(\varepsilon^{-\delta})
 \qquad(\varepsilon\downarrow0).
\]
\end{proposition}

\begin{proof}
Die $n$-te untere Lücke ist nach Skalierung mit $r_n^{-1}$ die Packung zum
Parameter $\theta_n$. Die uniforme Normierung liefert deshalb für $s>\delta$
\[
 \sum_{n\ge0}\sum_{C\in\mathcal A_n^{\mathrm G}}
 \operatorname{rad}(C)^s
 \ll_s\sum_{n\ge0}r_n^s<\infty.
\]
Die einzelne Lücke $n=0$ erzwingt die Divergenz für $s\le\delta$.
Entsprechend folgt für die Zählung die obere Schranke aus
\[
 \sum_{n\ge0}N_{\theta_n}(\varepsilon/r_n)
 \ll \varepsilon^{-\delta}\sum_{n\ge0}r_n^\delta,
\]
während $n=0$ die untere Schranke liefert. Die Cayley-Normierung der
Außenkreisfamilie besitzt auf jedem Wurzelintervall denselben gleichmäßigen
Kreisvergleich; die Spiegelung fügt nur weitere gleichartige Familien hinzu.
\end{proof}

\begin{remark}
Damit sind kritischer Exponent und Zählordnung nun für die gesamte
Halbkreispackung bestimmt. Die Zählhauptkonstante wird im folgenden
Unterabschnitt ohne effektiven Restterm bestimmt. Offen bleibt danach nur
noch die meromorphe Zeta-Fortsetzung mit kontrolliertem Fehler.
\end{remark}

\subsection{Reduktion der globalen Hauptkonstante}

Die Wurzelparameter der unteren Lücken besitzen eine besonders einfache
Grenzstruktur. Aus den Formeln für $\cosh$ folgt
\begin{equation}\label{eq:root-parameter-exact}
 r_n=\frac{2q^{-2n}}{(1+q^{-2n})^2},
 \qquad
 \theta_n
 =q\frac{1+q^{-2n-2}}{1+q^{-2n}}.
\end{equation}
Insbesondere konvergiert $\theta_n$ exponentiell gegen $q$. Sei
$\mathcal A_{n,\sigma}^{\sharp}$, $\sigma\in\{\mathrm G,\mathrm K\}$,
die nach einer Translation und Skalierung mit $r_n^{-1}$ normierte lokale
Packung in der $n$-ten rechten Lücke. Schreibe
$c_{n,\sigma}$ für ihre lokale Zählkonstante. Für die Geradenlücke ist
$c_{n,\mathrm G}=c_{\mathrm G}(\theta_n)$.

Die uniforme Normierung liefert die für die Summation nötige
Beschränktheit direkt. Für die untere Familie gilt wegen
Proposition~8.2 mit Konstanten $m,M>0$
\[
 N_{\mathcal G_2}(\varepsilon/m)
 \le N_{\mathcal G_\theta}(\varepsilon)
 \le N_{\mathcal G_2}(\varepsilon/M).
\]
Nach Multiplikation mit $\varepsilon^\delta$ und Grenzübergang folgt
\[
 c_{\mathrm G}(2)m^\delta
 \le c_{\mathrm G}(\theta)
 \le c_{\mathrm G}(2)M^\delta.
\]
Der Cayley-Vergleich liefert dieselbe Beschränktheit für die Außenfamilie.
Die analytische Abhängigkeit der Apollonischen Hauptkonstante von den
Anfangskreisen ist darüber hinaus in
\href{https://arxiv.org/abs/1702.02854}{Kesseböhmer--Kombrink}
beschrieben. Damit ist
\begin{equation}\label{eq:global-leading-constant}
 C_{\mathrm{glob}}
 =2\sum_{n\ge0}r_n^\delta
   \bigl(c_{n,\mathrm G}+c_{n,\mathrm K}\bigr)
\end{equation}
wohldefiniert und positiv. Der Faktor $2$ berücksichtigt die linke
Spiegelung. Formel \eqref{eq:root-parameter-exact} macht die untere
Teilreihe sogar rasch numerisch auswertbar:
\[
 \sum_{n\ge0}r_n^\delta c_{\mathrm G}(\theta_n)
 =c_{\mathrm G}(q)\sum_{n\ge0}r_n^\delta
  +\sum_{n\ge0}r_n^\delta
   \bigl(c_{\mathrm G}(\theta_n)-c_{\mathrm G}(q)\bigr).
\]

\begin{proposition}[Unbedingte globale Hauptasymptotik]
Für die vollständige Halbkreispackung gilt
\begin{equation}\label{eq:global-main-asymptotic}
 N_{\mathcal P}(\varepsilon)
 \sim C_{\mathrm{glob}}\varepsilon^{-\delta}
 \qquad(\varepsilon\downarrow0).
\end{equation}
Insbesondere besitzt die Zeta-Funktion die Abel-Grenze
\[
 \lim_{s\downarrow\delta}
 (s-\delta)Z_{\mathcal P}(s)=\delta C_{\mathrm{glob}}.
\]
\end{proposition}

\begin{proof}
Die uniformen Möbius- und Cayley-Normierungen liefern eine Konstante
$B>0$, so dass für alle $n$, alle $\sigma\in\{\mathrm G,\mathrm K\}$ und
alle $\varepsilon>0$
\[
 N_{\mathcal A_{n,\sigma}}(\varepsilon)
 \le B\left(\frac{r_n}{\varepsilon}\right)^\delta.
\]
Für ein festes $N$ liefert der lokale Kreiszählsatz nach Multiplikation mit
$\varepsilon^\delta$
\[
\lim_{\varepsilon\downarrow0}
\varepsilon^\delta
 2\sum_{n=0}^{N}\sum_{\sigma\in\{\mathrm G,\mathrm K\}}
 N_{\mathcal A_{n,\sigma}}(\varepsilon)
 =2\sum_{n=0}^{N}r_n^\delta
  \bigl(c_{n,\mathrm G}+c_{n,\mathrm K}\bigr).
\]
Andererseits ist der Rest gleichmäßig kontrolliert:
\[
 0\le 2\varepsilon^\delta
 \sum_{n>N}\sum_{\sigma\in\{\mathrm G,\mathrm K\}}
 N_{\mathcal A_{n,\sigma}}(\varepsilon)
 \le4B\sum_{n>N}r_n^\delta.
\]
Die rechte Seite verschwindet für $N\to\infty$. Dasselbe gilt für den Rest
der Reihe \eqref{eq:global-leading-constant}; die Skelettkette ist nach
Multiplikation mit $\varepsilon^\delta$ vernachlässigbar. Dies beweist
\eqref{eq:global-main-asymptotic}. Die Abel-Grenze folgt aus der
Mellin-Identität für $Z_{\mathcal P}$ und der Hauptasymptotik.
\end{proof}

\paragraph{Effektive Resthypothese.}
Es gebe ein $\eta\in(0,\delta)$, so dass für
$\sigma\in\{\mathrm G,\mathrm K\}$
\[
 N_{\mathcal A_{n,\sigma}^{\sharp}}(t)
 =c_{n,\sigma}t^{-\delta}
  +O\bigl(t^{-(\delta-\eta)}\bigr)
 \qquad(t\downarrow0)
\]
gilt, wobei die Konstante im $O$ unabhängig von $n$ und $\sigma$ ist.
Effektive Kreiszählsätze für feste Kleinsche Packungen sind bekannt
\href{https://arxiv.org/abs/2205.13004}{(Kontorovich--Lutsko)};
die hier noch zu prüfende Zusatzarbeit ist genau ihre uniforme Anwendung auf
unserem kompakten Formraum.

\begin{proposition}[Konditionale effektive Verfeinerung]
Unter der effektiven Resthypothese verbessert sich
\eqref{eq:global-main-asymptotic} zu
\begin{equation}\label{eq:global-exact-count}
 N_{\mathcal P}(\varepsilon)
 =C_{\mathrm{glob}}\varepsilon^{-\delta}
  +O\bigl(\varepsilon^{-(\delta-\eta)}\bigr).
\end{equation}
Folglich setzt sich die Radiuszeta meromorph in die Halbebene
$\operatorname{Re}s>\delta-\eta$ fort und besitzt bei $s=\delta$ den
einfachen Pol
\[
 Z_{\mathcal P}(s)
 =\frac{\delta C_{\mathrm{glob}}}{s-\delta}
  +\text{einen in }\operatorname{Re}s>\delta-\eta
  \text{ holomorphen Term}.
\]
\end{proposition}

\begin{proof}
Für die Lücken mit $r_n\gg\varepsilon$ wird die Resthypothese nach
Skalierung zu
\[
 N_{\mathcal A_{n,\sigma}}(\varepsilon)
 =c_{n,\sigma}r_n^\delta\varepsilon^{-\delta}
  +O\bigl(r_n^{\delta-\eta}
  \varepsilon^{-(\delta-\eta)}\bigr).
\]
Die Fehler summieren sich, weil
$\sum_n r_n^{\delta-\eta}<\infty$. Das Abschneiden der Hauptreihe bei
$r_n\asymp\varepsilon$ kostet nur $O(1)$, da
$\sum_{r_n\ll\varepsilon}r_n^\delta=O(\varepsilon^\delta)$.
Die Skelettkette trägt lediglich $O(\log(1/\varepsilon))$ Kreise bei.
Damit folgt \eqref{eq:global-exact-count}. Die Mellin-Identität
\[
 Z_{\mathcal P}(s)
 =s\int_0^{R}N_{\mathcal P}(\varepsilon)
  \varepsilon^{s-1}\,d\varepsilon
\]
liefert die Fortsetzung und den angegebenen Residuenwert.
\end{proof}

\section{Das projizierte Randmaß}

Sei $\beta$ das faire Bernoulli-Maß auf $\{L,R\}^{\mathbb N}$. Die
Stern--Brocot-Randkodierung schiebt $\beta$ auf das Minkowski-
Fragezeichenmaß $\mu_{?}$ auf $[0,1]$. Für jede feste geometrische
Wurzellücke ist die geometrische Kodierung eine projektive, auf dem
abgeschlossenen Grundintervall glatte Koordinatenänderung $h_n$. Also gilt
für das geometrische Randmaß
\[
 \nu_n=(h_n)_*\mu_{?}.
\]

\begin{proposition}[Singularität und typische Dimension]
Das faire geometrische Randmaß $\nu_n$ ist singulär-stetig und hat vollen
Träger auf seinem Grundintervall. Es ist exakt dimensional, und seine
Maßdimension ist die des Minkowski-Maßes. Sie erfüllt die rigorosen Schranken
\[
 0.874716305108207<\dim \nu_n<0.874716305108213.
\]
Insbesondere ist die zugehörige Lyapunov-Rate
\[
 \chi=\frac{\log2}{\dim\mu_{?}}
 \approx0.7924251286.
\]
\end{proposition}

\begin{proof}
Glatte projektive Koordinatenwechsel mit von Null verschiedener Ableitung
erhalten Singularität, Atomlosigkeit, vollen Träger und exakte Dimension.
Für das Minkowski-Maß sind diese Eigenschaften und rigorose Schranken für
die angegebene Dimension in
\href{https://arxiv.org/abs/0706.0453}{Kesseböhmer--Stratmann} und
\href{https://publicatio.bibl.u-szeged.hu/17247/1/corr3.pdf}{Mantica--M"{a}ki--Totik}
dokumentiert.
\end{proof}

Die Stern--Brocot-Multifraktaltheorie liefert zudem den passenden Rahmen für
die periodischen Adressen. Um die Dichte periodischer Dimensionswerte im
vollständigen Spektrum als eigenen Satz zu zitieren, bleibt eine saubere
Übertragung der parabolischen induzierten Kodierung auf unsere gewählte
Definition der lokalen Dimension zu schreiben. Die Frage ist damit nicht
mehr empirisch, sondern eine präzise Anwendung dieser Theorie.

\section{Der minimale generatorstabile Beobachtungsraum}

Die letzte offene Frage hat eine kanonische abstrakte Antwort. Sei $X$ der
Raum markierter Descartes-Konfigurationen, $\mathcal A^*$ das Monoid der
Generatorwörter und $\pi:X\to Y$ eine gewählte Beobachtung, etwa die
Projektion auf $[0,1]$. Definiere
\[
 x\sim_\pi y
 \quad\Longleftrightarrow\quad
 \pi(gx)=\pi(gy)\quad\text{für alle }g\in\mathcal A^*.
\]

\begin{proposition}[Kanonischer Beobachtungsquotient]
Jeder Apollonische Generator steigt zu einer wohldefinierten Funktion auf
$X/{\sim_\pi}$ ab. Dieser Quotient ist der kleinste generatorstabile Raum,
durch den $\pi$ faktorisiert: Jeder andere generatorstabile Zustandsraum,
der alle zukünftigen Beobachtungen bestimmt, faktorisiert auf
$X/{\sim_\pi}$.
\end{proposition}

\begin{proof}
Aus $x\sim_\pi y$ folgt für jedes Generatorwort $h$ und jeden Generator
$S_i$ die Gleichheit $\pi(hS_ix)=\pi(hS_iy)$. Also ist
$S_ix\sim_\pi S_iy$, und $S_i$ steigt ab. Jede weitere ausreichende
Beobachtung darf nur Paare identifizieren, die unter allen zukünftigen
Generatorwörtern dieselbe Projektion besitzen; daher faktorisiert sie über
den Quotienten.
\end{proof}

Dieser Satz löst die Minimalitätsfrage universell. Ob dieser Quotient in
unserer speziellen Geometrie durch wenige reelle Faserkoordinaten, etwa eine
markierte Descartes-Quadruple, beschrieben werden kann, bleibt eine echte
konkrete Frage.

\section{Lokale Dimensionen und die parabolische Kusp}

Nach Proposition~9.1 ist das faire Bernoulli-Maß auf jedem geometrischen
Randbaum ein glattes projektives Bild des Minkowski-Maßes. Insbesondere ist
die typische lokale Dimension kein bloßer numerischer Hinweis mehr.

\begin{proposition}[Lokale Dimension periodischer Adressen]
Sei $q$ ein periodisches Wort, dessen Matrixprodukt $M_q$ strikt positive
Einträge besitzt, und sei
$\rho_q>1$ dessen dominanter Eigenwert, dann gilt für die lokale Dimension
der periodischen Adresse
\begin{equation}\label{eq:periodic-dimension}
 d(q^\infty)=\frac{|q|\log2}{2\log\rho_q}.
\end{equation}
\end{proposition}
\begin{proof}
Da $M_q$ positiv ist, sind beide Komponenten von
$M_q^k(a,b)^{\mathsf T}$ für jeden positiven Startvektor $(a,b)$ zu
$\rho_q^k$ vergleichbar. Die geometrische Zylinderlänge ist daher nach
\eqref{eq:comparison} zu $\rho_q^{-2k}$ vergleichbar, während das faire
Bernoulli-Maß des Zylinders $2^{-k|q|}$ beträgt. Der Quotient der beiden
Logarithmen liefert \eqref{eq:periodic-dimension}.
\end{proof}

Zwei Beispiele sind
\[
 d((LR)^\infty)=\frac{\log2}{2\log\varphi}\approx0.72021,
\]
und
\[
 d((LR^2L)^\infty)=
 \frac{2\log2}{\log(3+2\sqrt2)}\approx0.78644.
\]
Die fast-parabolischen Wörter $q_n=L^{n-1}R$ liefern formal
$d(q_n^\infty)\to+\infty$. Damit wird die Kusp nicht als einzelner
Ausnahmeeffekt, sondern als Grenzbereich der periodischen Dynamik sichtbar.

\begin{proposition}[Dichte periodischer Dimensionswerte]
Die Menge der Dimensionswerte periodischer Adressen enthält eine in
\[
 \left[d((LR)^\infty),+\infty\right).
\]
dichte Teilmenge.
\end{proposition}

\begin{proof}
Setze $U=LR$ und $V_m=L^{m-1}R$ für $m\ge2$. Beide Matrizen sind positiv.
Für positive Matrizen $A,B$ folgt aus der Perron--Frobenius-Zerlegung
\[
 \log\rho(A^aB^b)
 =a\log\rho(A)+b\log\rho(B)+O(1)
\]
für $a,b\to\infty$. Wendet man \eqref{eq:periodic-dimension} auf die
Wörter $U^aV_m^b$ an, so nähern ihre Kehrwerte der Dimensionen genau die
gewichteten Mittelwerte von $d(U^\infty)^{-1}$ und
$d(V_m^\infty)^{-1}$ an. Mit dichten rationalen Verhältnissen $a:b$
werden damit alle Werte im Intervall
\[
 [d(U^\infty),d(V_m^\infty)]
\]
dicht getroffen. Schließlich gilt
\[
 V_m=\begin{pmatrix}1&1\\m-1&m\end{pmatrix},
 \qquad
 \rho(V_m)
 =\frac{m+1+\sqrt{(m+1)^2-4}}2,
\]
also $d(V_m^\infty)\to+\infty$. Die Vereinigung dieser Intervalle liefert
die Behauptung.
\end{proof}

\begin{proposition}[Scharfer linker Spektrumsrand]
Sei $x$ eine Stern--Brocot-Adresse, für die die lokale Dimension existiert.
Dann gilt
\[
 d(x)\ge \frac{\log2}{2\log\varphi}.
\]
Gleichheit wird von der alternierenden periodischen Adresse
$(LR)^\infty$ erreicht. Damit ist
\[
 \frac{\log2}{2\log\varphi}
\]
der linke Endpunkt des vollständigen lokalen Dimensionsspektrums.
\end{proposition}

\begin{proof}
Seien $F_0=0$, $F_1=1$ die Fibonacci-Zahlen. Für ein Wort $w$ der Länge
$n$ schreibe $(u_w,v_w)=M_w(1,1)^{\mathsf T}$. Induktion über $n$ liefert
\[
 \max(u_w,v_w)\le F_{n+2},
 \qquad
 u_w+v_w\le F_{n+3}.
\]
Denn eine Anwendung von $L$ oder $R$ ersetzt eines der beiden Paare durch
$u_w+v_w$; zugleich wächst die Summe höchstens um die bisher größere
Komponente. Aus \eqref{eq:comparison} folgt daher für die
Zylinderlänge
\[
 \ell_w\ge \frac{c}{F_{n+2}^2}.
\]
Da $F_{n+2}=\varphi^{n+O(1)}$ und der faire Zylinder die Masse $2^{-n}$
hat, ist jeder existierende Dimensionsgrenzwert mindestens
$\log2/(2\log\varphi)$. Die Gleichheit für $(LR)^\infty$ folgt aus
\eqref{eq:periodic-dimension}.
\end{proof}

\begin{corollary}[Abschluss des lokalen Spektrums]
Für jeden festen geometrischen Randbaum sei
\[
 \operatorname{Spec}_{\mathrm{loc}}
 =\{d(x):\text{ die lokale Dimension von }x\text{ existiert}\}.
\]
Dann gilt im erweiterten Intervall
\[
 \overline{\operatorname{Spec}_{\mathrm{loc}}}
 =\left[\frac{\log2}{2\log\varphi},+\infty\right].
\]
\end{corollary}

\begin{proof}
Proposition~11.3 gibt die linke Einschlussrichtung. Proposition~11.2
liefert eine dichte Teilmenge des ganzen rechten Halbstrahls durch
periodische Adressen; die fast-parabolischen periodischen Wörter lassen
deren Dimensionen gegen $+\infty$ gehen.
\end{proof}

\begin{corollary}[Wurzelunabhängigkeit]
Die lokalen Dimensionsspektren und die Hausdorff-Dimensionen ihrer
Niveaumengen sind für alle geometrischen Wurzellücken $G_n$ gleich.
\end{corollary}

\begin{proof}
Nach Proposition~9.1 ist jedes geometrische Randmaß ein Bild des
Minkowski-Maßes unter einer glatten projektiven Koordinatenänderung mit
von Null verschiedener Ableitung auf dem kompakten Grundintervall.
Solche Abbildungen und ihre Inversen sind bi-Lipschitz. Sie erhalten
lokale Dimensionen und Hausdorff-Dimensionen von Niveaumengen.
\end{proof}

\begin{proposition}[Nichtperiodische Minimierer]
Setze
\[
 d_*:=\frac{\log2}{2\log\varphi},
 \qquad E_*:=\{x:d(x)=d_*\}.
\]
Dann enthält $E_*$ eine überabzählbare Menge nichtperiodischer Adressen.
Diese konstruierte Teilmenge besitzt Hausdorff-Dimension $0$.
\end{proposition}

\begin{proof}
Setze $U=LR$, $m_j=2^j$ und wähle für jede Folge
$\varepsilon=(\varepsilon_j)_{j\ge1}\in\{0,1\}^{\mathbb N}$ die Adresse
\[
 w(\varepsilon)=U^{m_1}B_{\varepsilon_1}
 U^{m_2}B_{\varepsilon_2}\cdots,
 \qquad B_0=LL,\quad B_1=RR.
\]
Der Perronwert von $U$ ist $\varphi^2$. Daher existiert $A>0$ mit
\[
 \|U^m z\|_1\ge A\varphi^{2m}\|z\|_1
 \qquad(m\ge1,\ z\in\R_{>0}^2).
\]
Die beiden Blöcke $B_0,B_1$ verkleinern die $1$-Norm nicht. Nach den
ersten $J$ Blöcken hat das Wort die Länge
\[
 N_J=2\sum_{j=1}^Jm_j+2J,
\]
und seine Parameter erfüllen
\[
 \log(u_{w}v_{w})=2N_J\log\varphi+O(J).
\]
Tatsächlich gibt die Fibonacci-Schranke die entgegengesetzte obere
Abschätzung für die $1$-Norm. Nach einem $U$-Block und nach einem der
beiden Blöcke $B_i$ sind die Komponenten außerdem in einem festen
Verhältnis, sodass ihr Produkt mit dem Quadrat der $1$-Norm vergleichbar
ist. Dies gilt ebenso für Präfixe innerhalb eines $U$-Blocks, denn
$J=O(\log N_J)$. Der geometrische Längenvergleich liefert folglich
\[
 -\log\ell_{n,w_1\cdots w_N}
 =2N\log\varphi+o(N).
\]
Da der faire Zylinder die Masse $2^{-N}$ hat, ist die lokale Dimension
gleich $d_*$. Die Blöcke $U^{m_j}$ enthalten unbeschränkt lange
alternierende Läufe, während nach jedem solchen Lauf ein Doppelbuchstabe
steht. Die Adressen sind daher verschieden und nichtperiodisch.

Auf Stufe $J$ überdecken $2^J$ Zylinder die konstruierte Menge; ihre
Längen sind höchstens
$\exp(-2N_J\log\varphi+O(J))$. Für jedes $s>0$ geht die zugehörige
$s$-Hausdorffsumme gegen $0$. Die Teilmenge hat somit
Hausdorff-Dimension $0$.
\end{proof}

\begin{proposition}[Aperiodische Realisierung aller endlichen Werte]
Für jedes
\[
 d\in\left[\frac{\log2}{2\log\varphi},+\infty\right)
\]
gibt es eine nichtperiodische Adresse, deren lokale Dimension gleich $d$
ist.
\end{proposition}

\begin{proof}
Der linke Endpunkt ist durch Proposition~11.6 realisiert. Sei nun
$d>d_*$ und setze $\gamma_d=\log2/(2d)$. Für ein positives Wort $W$
schreibe
\[
 \gamma(W)=\frac{\log\rho_W}{|W|}.
\]
Mit $U=LR$ gilt $\gamma(U)=\log\varphi$. Für
$V_m=L^{m-1}R$ gilt dagegen $\gamma(V_m)\to0$. Wähle also $m$ mit
\[
 \gamma(V_m)<\gamma_d<\gamma(U)
\]
und $\theta\in(0,1)$ so, dass
\[
 \gamma_d=\theta\gamma(U)+(1-\theta)\gamma(V_m).
\]

Für $j\ge1$ wähle ganze Zahlen $p_j,q_j\ge1$ mit
\[
 2p_j+m q_j=j+O(1),
 \qquad
 \frac{2p_j}{2p_j+m q_j}\longrightarrow\theta,
\]
und betrachte die Adresse
\[
 U^{p_1}V_m^{q_1}U^{p_2}V_m^{q_2}\cdots.
\]
Die Perron--Frobenius-Zerlegung der beiden positiven Blöcke liefert nach
$J$ Makroblöcken
\[
 \log\|M_w\|_1
 =\sum_{j\le J}\bigl(p_j\log\rho_U+q_j\log\rho_{V_m}\bigr)+O(J).
\]
Die Wortlänge ist von Ordnung $J^2$, während ein Makroblock nur Länge
$O(J)$ besitzt. Dieselbe Grenzrate gilt folglich für alle Präfixe. Der
geometrische Längenvergleich ergibt
\[
 -\log\ell_{n,w_1\cdots w_N}=2\gamma_dN+o(N),
\]
also lokale Dimension $d$. Die Blöcke $U^{p_j}$ enthalten unbeschränkt
lange alternierende Läufe, während die Blöcke $V_m^{q_j}$ unbeschränkt
oft Doppelbuchstaben erzeugen. Die Adresse kann daher nicht periodisch
sein.
\end{proof}

Die typische lokale Dimension ist die Maßdimension des Minkowski-Maßes und
erfüllt somit
\[
 0.874716305108207<d_{\mathrm{typ}}<0.874716305108213,
 \qquad \chi\approx0.7924251286.
\]
Sie entspricht dem fast sicheren Grenzwert
\[
 \chi=\lim_{n\to\infty}-\frac1n\log\ell_{w_1\cdots w_n}
\]
in der Stern--Brocot-Kodierung.

\begin{proposition}[Kuspuntergrenze]
Für die Menge $E(+\infty)$ der Randpunkte mit unendlicher lokaler
Dimension gilt
\[
 \Haus E(+\infty)\ge\frac12.
\]
\end{proposition}

\begin{proof}
Setze $N_j=(j+2)^2$ und wähle in der Stufe $j$ die geraden Schrittweiten
\[
 D_j=\{N_j,N_j+2,\ldots,2N_j\}.
\]
Wir betrachten alle unendlichen Konkatenationen der Blöcke $L^aR$ mit
$a\in D_j$ in Stufe $j$. Nach einem ganzen Block sind die beiden
Parameter in einem festen Verhältnis. Die direkte Rekursion zeigt daher
für einen Zylinder der Stufe $J$
\[
 \ell\asymp\prod_{j=1}^J a_j^{-2}.
\]
Die Auswahl jedes zweiten Werts lässt zwischen zwei gewählten Kindern
jeweils mindestens einen vollständigen Zwischenzylinder. Die
Separation ist somit mit der Zylinderlänge vergleichbar. Das
Moran-Massenprinzip liefert deshalb
\[
 \Haus\mathcal C
 =\lim_{J\to\infty}
 \frac{\sum_{j\le J}\log|D_j|}
 {2\sum_{j\le J}\log N_j}=\frac12
\]
für die konstruierte Menge $\mathcal C$.

Die Worttiefe eines Stufe-$J$-Zylinders ist von Ordnung
$\sum_{j\le J}N_j\asymp J^3$, während sein negativer Längenlogarithmus
von Ordnung $\sum_{j\le J}\log N_j=O(J\log J)$ ist. Innerhalb eines
Blocks fügt ein Präfix $L^t$ nur einen zusätzlichen Term
$O(\log(t+1))$ zum Längenlogarithmus hinzu. Folglich ist für jedes
Element von $\mathcal C$
\[
 \frac{n\log2}{-\log\ell_{w_1\cdots w_n}}\longrightarrow+\infty.
\]
Also ist $\mathcal C\subseteq E(+\infty)$, was die Behauptung beweist.
\end{proof}

\begin{conjecture}
Es gilt die scharfe Gleichheit
\[
 \Haus E(+\infty)=\frac12.
\]
Die noch fehlende Richtung ist eine obere Schranke für alle Adressen mit
subexponentiellem geometrischem Zylinderschwund.
\end{conjecture}

\subsection{Nächster Forschungsblock: Tiefenstatistik}

Für eine feste Wurzellücke $G_n$ und eine Baumtiefe $k$ definieren wir
die Ebenensumme
\[
 \mathcal Z_{n,k}(s)=\sum_{|w|=k}\ell_{n,w}^{s}.
\]
Sie misst, wie sich die geometrische Länge auf die $2^k$ Lücken einer
symbolischen Tiefe verteilt.

\begin{proposition}[Tiefenerhaltung und Extremalskalen]
Für jedes $n$ und $k\ge0$ gilt
\[
 \mathcal Z_{n,k}(1)=\ell_{n,\varnothing}
 =\frac{2}{a_nb_n}.
\]
Außerdem existieren positive, von $k$ unabhängige Konstanten, so dass
\[
 \max_{|w|=k}\ell_{n,w}\asymp k^{-1},
 \qquad
 \min_{|w|=k}\ell_{n,w}\asymp\varphi^{-2k}.
\]
\end{proposition}

\begin{proof}
Die beiden Kinder eines Parameterpaars $(a,b)$ haben nach
\eqref{eq:gap-length-exact} Längen
\[
 \frac{2}{a(a+b)}\quad\text{und}\quad
 \frac{2}{b(a+b)},
\]
deren Summe $2/(ab)$ ist. Iteration beweist die erste Behauptung.

Für den oberen Extremalwert wächst das Produkt der Parameter bei jedem
Schritt mindestens um $\min(a_n,b_n)^2$:
\[
 a(a+b)=ab+a^2,\qquad (a+b)b=ab+b^2.
\]
Daher ist jede Lücke der Tiefe $k$ höchstens von Ordnung $k^{-1}$.
Der parabolische Ast $L^k$ besitzt dagegen das Produkt
$a_n(b_n+ka_n)$ und damit Länge von Ordnung $k^{-1}$.

Die Fibonacci-Schranke aus Proposition~11.3 und
\eqref{eq:geometric-comparison} liefert die untere Schranke
$\ell_{n,w}\gg\varphi^{-2k}$ für alle Wörter der Tiefe $k$. Der
alternierende Ast $U^{\lfloor k/2\rfloor}$, bei ungeradem $k$ um einen
Buchstaben ergänzt, liefert die umgekehrte Schranke.
\end{proof}

\begin{corollary}[Erster Rahmen für die Ebenensummen]
Für jedes $s>1$ gibt es Konstanten $c_{n,s},C_{n,s}>0$ mit
\[
 c_{n,s}k^{-s}\le\mathcal Z_{n,k}(s)
 \le C_{n,s}k^{1-s}
 \qquad(k\ge1).
\]
\end{corollary}

\begin{proof}
Die untere Schranke folgt vom parabolischen Ast. Für die obere gilt
\[
 \mathcal Z_{n,k}(s)
 \le \left(\max_{|w|=k}\ell_{n,w}\right)^{s-1}
 \sum_{|w|=k}\ell_{n,w},
\]
und die vorangehende Proposition liefert die Behauptung.
\end{proof}

Damit ist die vermutete scharfe Form
\[
 \mathcal Z_{n,k}(s)\sim C_{n,s}k^{-s}\qquad(s>1)
\]
auf eine konkrete Verteilungsfrage reduziert: Es muss kontrolliert werden,
wie viele Lücken einer Tiefe mit dem parabolischen Maßstab $k^{-1}$
vergleichbar bleiben.

\subsubsection{Parabolische Randstrahlen im Referenzmodell}

Für das Referenzintervall mit Parameterpaar $(1,1)$ sei
\[
 \mathcal Z_k^{\mathrm{ref}}(s)
 =\sum_{|w|=k}\left(\frac{2}{u_wv_w}\right)^s.
\]
Die maximalen Endläufe liefern bereits eine explizite untere
Asymptotik.

\begin{proposition}[Parabolischer Liminf]
Für $s>1$ gilt
\[
 \liminf_{k\to\infty}k^s\mathcal Z_k^{\mathrm{ref}}(s)
 \ge 2^{s+1}\frac{\zeta(2s-1)}{\zeta(2s)}.
\]
\end{proposition}

\begin{proof}
Betrachte zunächst Wörter, die nach einem festen Präfix $v$ nur noch mit
$L$ fortgesetzt werden. Endet $v$ mit $R$ oder ist $v$ leer und besitzt
es das Paar $(a,b)$, so hat $vL^m$ die Länge
\[
 \frac{2}{a(b+ma)}.
\]
Für $k=|v|+m$ trägt dieser Strahl daher im Limes
$(2/a^2)^s$ zu $k^s\mathcal Z_k^{\mathrm{ref}}(s)$ bei. Die möglichen
Paare sind genau $(1,1)$ und die teilerfremden Paare $a>b\ge1$. Ihre
Summe ist
\[
 \sum_{a\ge1}\frac{\varphi(a)}{a^{2s}}
 =\frac{\zeta(2s-1)}{\zeta(2s)}.
\]
Die spiegelbildlichen $R$-Strahlen liefern denselben Beitrag. Jede
endliche Teilmenge dieser Strahlen ist ab einer festen Tiefe disjunkt.
Zuerst für endliche Teilmengen den Grenzwert zu bilden und dann zu
erschöpfen beweist die Behauptung.
\end{proof}

\subsubsection{Die Zwei-Wechsel-Familie}

Sei $\mathcal W_{k,2}$ die Menge der Wörter der Länge $k$ mit genau
zwei Symbolwechseln und setze
\[
 \mathcal Z_{k,2}^{\mathrm{ref}}(s)
 =\sum_{w\in\mathcal W_{k,2}}
 \left(\frac{2}{u_wv_w}\right)^s.
\]
Diese erste nichttriviale Wechselklasse bestätigt die parabolische
Kandidatenkonstante bereits exakt.

\begin{proposition}[Zwei-Wechsel-Beitrag]
Für jedes $s>1$ gilt
\[
 \lim_{k\to\infty}
 k^s\mathcal Z_{k,2}^{\mathrm{ref}}(s)
 =
 2^{s+1}\sum_{r,j\ge1}
 \frac{1}{\bigl(1+j(r+1)\bigr)^{2s}}.
\]
\end{proposition}

\begin{proof}
Die Wörter in $\mathcal W_{k,2}$ sind genau
\[
 L^rR^jL^m
 \quad\text{und}\quad
 R^rL^jR^m,
 \qquad r,j,m\ge1,\qquad r+j+m=k.
\]
Nach $L^rR^j$ lautet das Parameterpaar
\[
 (a,b)=\bigl(1+j(r+1),\,r+1\bigr).
\]
Der letzte $L$-Lauf liefert daher die Länge
\[
 \frac{2}{a(b+ma)}.
\]
Die spiegelbildliche Familie hat dieselbe Länge. Also ist
\[
 k^s\mathcal Z_{k,2}^{\mathrm{ref}}(s)
 =
 2^{s+1}
 \sum_{\substack{r,j,m\ge1\\r+j+m=k}}
 \left(\frac{k}{a(b+ma)}\right)^s.
\]

Zunächst beschränken wir die Summe auf $m\ge k/2$. Dann gilt
\[
 \frac{k}{a(b+ma)}
 \le\frac{k}{ma^2}
 \le\frac{2}{a^2}.
\]
Die rechte Seite ist über $r,j$ summierbar, denn
$a=1+j(r+1)\ge j(r+1)$ und
\[
 \sum_{r,j\ge1}a^{-2s}<\infty.
\]
Für festes $r,j$ konvergiert der Summand gegen $a^{-2s}$.
Dominierte Konvergenz liefert somit den angegebenen Grenzwert für
den Bereich $m\ge k/2$.

Im Restbereich $m<k/2$ gilt mit $r+j=k-m>k/2$
\[
 a=1+j(r+1)\ge r+j+1>\frac{k}{2}.
\]
Folglich ist
\[
 \sum_{\substack{r,j\ge1\\r+j=k-m}}a^{-2s}
 \le k\left(\frac{2}{k}\right)^{2s}.
\]
Der gesamte Rest ist daher durch
\[
 2^{3s+1}k^{1-s}
 \sum_{m<k/2}m^{-s}
\]
beschränkt und verschwindet für $k\to\infty$, weil $s>1$.
\end{proof}

\begin{remark}
Die zwei Wechsel erzwingen die Kongruenzform
\[
 a=1+jb,\qquad b=r+1\ge2.
\]
Damit erscheint nur eine spezielle Teilmenge der teilerfremden Paare.
Bei beliebig vielen Wechseln entstehen dagegen alle Paare
$(a,b)$ mit $1\le b<a$ und $\gcd(a,b)=1$; ihre Anzahl für festes
$a$ ist $\varphi(a)$. Die Proposition macht somit sichtbar, wie die
volle Totientenkonstante schrittweise aus den Wechselklassen entsteht.
\end{remark}

\begin{conjecture}[Scharfe parabolische Ebenenasymptotik]
Für jedes $s>1$ gilt im Referenzmodell
\[
 \mathcal Z_k^{\mathrm{ref}}(s)
 \sim 2^{s+1}\frac{\zeta(2s-1)}{\zeta(2s)}\,k^{-s}.
\]
Die exakte Auswertung bis Tiefe $26$ stützt dies: Für $s=2$ ist
$k^2\mathcal Z_k^{\mathrm{ref}}(2)=8.46727$, während die
Kandidatenkonstante $8\zeta(3)/\zeta(4)$ ungefähr $8.88501$ beträgt.
\end{conjecture}

\begin{remark}[Numerische Signatur der Randstrahlen]
Die exakte Baumrechnung trennt die W\"orter auch nach der Zahl ihrer
Symbolwechsel. Bei Tiefe $26$ und $s=2$ entfallen $87{,}61\%$ der
Ebenensumme auf die beiden wechselarmen Randstrahlen ohne Wechsel,
$8{,}73\%$ auf W\"orter mit genau einem Wechsel und nur $3{,}66\%$ auf
alle \"ubrigen W\"orter. Auch nach der L\"ange des maximalen Endlaufs
konzentriert sich die Masse am Rand: Die Anteile der Endlaufl\"angen
$26,25,24,23$ sind jeweils ungef\"ahr
\[
 0.87612,\quad 0.06139,\quad 0.02628,\quad 0.01274.
\]
Dies ist nur Evidenz, isoliert aber genau den Mechanismus der folgenden
Erneuerungsfrage. Das auswertbare Skript
\texttt{tools/depth-stats.mjs} dokumentiert die Rechnung.
\end{remark}

\begin{remark}[Zeta-Konsistenz der Kandidatenkonstante]
Die Tiefensummen zerlegen die vollständige Referenzzeta zu
\[
 \sum_{k\ge0}\mathcal Z_k^{\mathrm{ref}}(s)
 =2^s\frac{\zeta(s)^2}{\zeta(2s)}.
\]
Für
\[
 C_s=2^{s+1}\frac{\zeta(2s-1)}{\zeta(2s)}
\]
gilt dagegen
\[
 C_s\zeta(s)\sim\frac{12}{\pi^2}(s-1)^{-2}
 \qquad(s\downarrow1).
\]
Dies stimmt exakt mit dem Doppelpolkoeffizienten der vollständigen
Referenzzeta überein. Die parabolische Ebenenasymptotik besitzt damit
die korrekte globale Zeta-Normierung.
\end{remark}

\subsubsection{Induzierte parabolische Erneuerung}

Für $j\ge1$ sei $\mathcal V_j^R$ die Menge der Wörter der Länge $j$,
die mit $R$ enden; außerdem setze
$\mathcal V_0^R=\{\varnothing\}$. Hat
$v\in\mathcal V_j^R$ das Parameterpaar $(a_v,b_v)$, so definieren wir
die einseitige Erneuerungsmasse
\[
 A_j(s)=\sum_{v\in\mathcal V_j^R}a_v^{-2s}.
\]

\begin{proposition}[Exakte Erneuerungsidentität]
Für jedes $s>0$ und jede Tiefe $k\ge1$ gilt
\[
 \mathcal Z_k^{\mathrm{ref}}(s)
 =
 2^{s+1}\sum_{j=0}^{k-1}
 \sum_{v\in\mathcal V_j^R}
 \frac{1}{\bigl(a_v(b_v+(k-j)a_v)\bigr)^s}.
\]
Ferner gilt für $s>1$
\[
 \sum_{j\ge0}A_j(s)
 =
 \sum_{a\ge1}\frac{\varphi(a)}{a^{2s}}
 =
 \frac{\zeta(2s-1)}{\zeta(2s)}.
\]
\end{proposition}

\begin{proof}
Jedes Wort der Tiefe $k$, das mit $L$ endet, besitzt eindeutig die Form
$vL^m$ mit $v\in\mathcal V_j^R$, $m=k-j\ge1$. Seine Länge ist
$2/[a_v(b_v+ma_v)]$. Die Wörter mit terminalem $R$ sind das Spiegelbild;
dies erklärt den Faktor $2^{s+1}$.

Die endlichen Wörter in
$\bigsqcup_{j\ge0}\mathcal V_j^R$ entsprechen bijektiv den Paaren
$(1,1)$ sowie den positiven teilerfremden Paaren $(a,b)$ mit $a>b$.
Für festes $a\ge2$ gibt es genau $\varphi(a)$ solche Werte von $b$;
der leere Präfix liefert den Term $a=1$.
\end{proof}

\begin{proposition}[Ausreichendes Kernkriterium]
Sei $s>1$. Gilt
\[
 A_j(s)\ll_s (j+1)^{-2s},
\]
so folgt die scharfe parabolische Ebenenasymptotik
\[
 \mathcal Z_k^{\mathrm{ref}}(s)
 \sim
 2^{s+1}\frac{\zeta(2s-1)}{\zeta(2s)}\,k^{-s}.
\]
\end{proposition}

\begin{proof}
Für festes $J$ liefert die Erneuerungsidentität
\[
 \lim_{k\to\infty}
 k^s\,2^{s+1}\sum_{j=0}^{J}
 \sum_{v\in\mathcal V_j^R}
 \frac{1}{\bigl(a_v(b_v+(k-j)a_v)\bigr)^s}
 =
 2^{s+1}\sum_{j=0}^{J}A_j(s).
\]
Bezeichne den Beitrag der Summanden mit $j>J$ in der
Erneuerungsidentität durch $\mathcal R_{k,J}(s)$. Aus
$b_v+(k-j)a_v\ge(k-j)a_v$ folgt
\[
 0\le k^s\mathcal R_{k,J}(s)
 \le
 2^{s+1}k^s
 \sum_{j=J+1}^{k-1}(k-j)^{-s}A_j(s).
\]
Im Bereich $j\le k/2$ ist $k/(k-j)\le2$, sodass dieser Ausdruck durch
ein konstantes Vielfaches von
$\sum_{j>J}(j+1)^{-2s}$ beschränkt ist. Im Bereich $j>k/2$ ist er
durch
\[
 C_s k^{-s}\sum_{m=1}^{\lfloor k/2\rfloor}m^{-s}
\]
beschränkt und verschwindet für $k\to\infty$. Zuerst
$k\to\infty$ und dann $J\to\infty$ liefert die Behauptung.
\end{proof}

\paragraph{Kernfrage.}
Die noch offene Aufgabe ist damit auf die einseitige Abschätzung
\[
 A_j(s)\ll_s j^{-2s}
 \qquad(s>1)
\]
reduziert. Sie ist scharf in der Größenordnung, denn der parabolische
Präfix $R^j$ liefert bereits den Summanden $(j+1)^{-2s}$. Im Unterschied
zur ursprünglichen Ebenensumme enthält $A_j(s)$ nur die festgehaltene
Koordinate vor dem terminalen Lauf. Diese einseitige Form ist der
natürliche Ansatzpunkt für eine induzierte parabolische Erneuerung.


\subsubsection{Induzierter parabolischer Transferoperator}

Die Kernmasse besitzt eine dynamische Darstellung, die die parabolischen
Läufe exakt aussondert. Für ein Stern--Brocot-Paar $(u,v)$ setze
\[
 a=\max(u,v),\qquad x=\frac{\min(u,v)}{\max(u,v)}\in(0,1].
\]
Nach jedem Schritt normieren wir wieder so, dass die größere Komponente
an erster Stelle steht. Die beiden möglichen Quotientenabbildungen sind
\[
 T_0(x)=\frac{x}{1+x},
 \qquad
 T_1(x)=\frac{1}{1+x};
\]
in beiden Fällen wird die größere Komponente mit $1+x$ multipliziert.

\begin{proposition}[Transferdarstellung der Kernmasse]
Definiere
\[
 B_k(s)=\sum_{|w|=k}\max(u_w,v_w)^{-2s}.
\]
Dann gilt für $k\ge1$
\[
 B_k(s)=2A_k(s).
\]
Mit dem gewichteten Transferoperator
\[
 (\mathcal L_sf)(x)
 =(1+x)^{-2s}
 \left(
 f\!\left(\frac{x}{1+x}\right)
 +f\!\left(\frac{1}{1+x}\right)
 \right)
\]
gilt ferner
\[
 B_k(s)=(\mathcal L_s^k\mathbf 1)(1).
\]
\end{proposition}

\begin{proof}
Aus einem normierten Paar $(a,b)$ mit $b=ax$ entstehen, nach erneuter
Normierung, die Paare
\[
 (a+b,b)
 \quad\text{und}\quad
 (a+b,a).
\]
Beide haben die größere Komponente $a(1+x)$ und die angegebenen
Quotienten. Die Iteration liefert die Transferdarstellung. Die Spiegelung
vertauscht die Wörter mit terminalem $L$ und terminalem $R$, ohne die
größere Komponente zu ändern; daher ist $B_k=2A_k$ für $k\ge1$.
\end{proof}

Die parabolische Verzweigung $T_0$ ist exakt integrierbar:
\[
 T_0^m(x)=\frac{x}{1+mx},
 \qquad
 \prod_{\ell=0}^{m-1}
 \bigl(1+T_0^\ell(x)\bigr)^{-2s}
 =(1+mx)^{-2s}.
\]
Induziert man auf dem zentralen Intervall
$Y=[1/2,1]$, so besitzt ein Rückkehrblock
$T_0^mT_1$ die Abbildung und das Gewicht
\[
 G_m(x)=\frac{1+mx}{1+(m+1)x},
 \qquad
 g_m(x)=\bigl(1+(m+1)x\bigr)^{-2s}.
\]
Damit lautet der induzierte Operator
\[
 (\mathcal G_sf)(x)
 =\sum_{m\ge0}g_m(x)f(G_m(x)),
 \qquad x\in Y.
\]
Für $s>1/2$ sind die Gewichte auf $Y$ gleichmäßig summierbar, denn
\[
 g_m(x)\ll_s (m+1)^{-2s}.
\]

\begin{proposition}[Uniforme Verzerrung und induzierte Resolvente]
Sei $\omega=(m_1,\ldots,m_n)$ ein endliches Wort in den induzierten
Zweigen und
\[
 G_\omega=G_{m_n}\circ\cdots\circ G_{m_1}.
\]
Dann gilt
\[
 \sup_{x,y\in Y}
 \frac{|G_\omega'(x)|}{|G_\omega'(y)|}\le4.
\]
Ferner gilt für jedes $s>1$
\[
 \|\mathcal G_s^n\|_{C(Y)\to C(Y)}
 \le
 4^s\left(\frac49\right)^{n(s-1)}.
\]
Insbesondere konvergiert
\[
 \sum_{n\ge0}\mathcal G_s^n
\]
in der Operatornorm auf $C(Y)$.
\end{proposition}

\begin{proof}
Jede Komposition $G_\omega$ ist eine Möbius-Abbildung mit nichtnegativen
Koeffizienten. Ihre Ableitung hat daher die Form
\[
 |G_\omega'(x)|=\frac{c}{(ax+b)^2}
 \qquad(a,b,c\ge0).
\]
Für $x,y\in[1/2,1]$ folgt
\[
 \frac{ax+b}{ay+b}\le2,
\]
woraus die Verzerrungsschranke folgt.

Die Intervalle
\[
 G_m(Y)=
 \left[
 \frac{m+1}{m+2},
 \frac{m+2}{m+3}
 \right]
\]
zerlegen $Y$ bis auf Endpunkte. Iteration zeigt, dass auch die
Zylinderintervalle $G_\omega(Y)$ einer festen Tiefe $n$ paarweise
disjunkt sind und zusammen $Y$ überdecken. Setze
\[
 H_n(x)=(\mathcal G_s^n\mathbf 1)(x)
 =\sum_{|\omega|=n}|G_\omega'(x)|^s.
\]
Da jede einzelne Ableitung höchstens $4/9$ beträgt, gilt
\[
 \int_YH_n(x)\,dx
 \le
 \left(\frac49\right)^{n(s-1)}
 \sum_{|\omega|=n}\int_Y|G_\omega'(x)|\,dx
 =
 \frac12\left(\frac49\right)^{n(s-1)}.
\]
Die Verzerrungsschranke überträgt diese mittlere Schranke auf jede Stelle
von $Y$ und liefert
\[
 \|H_n\|_\infty
 \le4^s\left(\frac49\right)^{n(s-1)}.
\]
Für $f\in C(Y)$ folgt
$|\mathcal G_s^nf|\le\|f\|_\infty H_n$, was die
Operatorabschätzung und die Resolventenaussage beweist.
\end{proof}

\begin{proposition}[Rang-eins-Schwanz der Rückkehroperatoren]
Für $m\ge0$ definiere auf $C^1(Y)$ den Rückkehroperator
\[
 (R_mf)(x)
 =
 \bigl(1+(m+1)x\bigr)^{-2s}f(G_m(x))
\]
und den Rang-eins-Operator
\[
 (\Pi_sf)(x)=x^{-2s}f(1).
\]
Für jedes kompakte Teilgebiet $K\subset\mathbb C$ gilt
gleichmäßig für $s\in K$
\[
 \left\|m^{2s}R_m(s)-\Pi_s\right\|_{C^1(Y)\to C^1(Y)}
 \ll_K\frac1m
 \qquad(m\to\infty).
\]
\end{proposition}

\begin{proof}
Für $x\in Y$ gilt gleichmäßig für $s$ in kompakten Teilgebieten von
$\mathbb C$
\[
 m^{2s}\bigl(1+(m+1)x\bigr)^{-2s}
 =
 x^{-2s}\left(1+O_K\!\left(\frac1m\right)\right),
\]
und dieselbe Abschätzung gilt nach Ableitung nach $x$. Ferner ist
\[
 1-G_m(x)
 =
 \frac{x}{1+(m+1)x}
 =O\!\left(\frac1m\right),
 \qquad
 G_m'(x)=O\!\left(\frac1{m^2}\right).
\]
Daher folgt für $f\in C^1(Y)$
\[
 \|f\circ G_m-f(1)\|_{C^1}
 \ll\frac{\|f\|_{C^1}}m.
\]
Multiplikation der beiden Abschätzungen liefert die Behauptung.
\end{proof}

\begin{proposition}[Starke induzierte Resolvente]
Für jedes $s>1$ konvergiert
\[
 \sum_{n\ge0}\mathcal G_s^n
\]
auch in der Operatornorm auf $C^1(Y)$.
\end{proposition}

\begin{proof}
Schreibe $g_m(x)=|G_m'(x)|^s$. Dann gilt
\[
 \|(\mathcal G_sf)'\|_\infty
 \le D_s\|f\|_\infty+\kappa_s\|f'\|_\infty,
\]
wobei
\[
 D_s=\sum_{m\ge0}\|g_m'\|_\infty<\infty,
 \qquad
 \kappa_s\le
 \sum_{m\ge0}\left(\frac2{m+3}\right)^{2s+2}<1.
\]
Zusammen mit der bereits bewiesenen exponentiellen Abschätzung auf
$C(Y)$ ergibt Iteration eine exponentielle Abschätzung auf $C^1(Y)$.
Damit folgt die Resolventenaussage.
\end{proof}

\begin{remark}[Defekter induzierter Bereich]
Für die induzierten Rückkehrgewichte auf $Y$ setze
\[
 q_s=\sum_{m\ge0}
 \left(\frac{2}{m+3}\right)^{2s}
 =2^{2s}\bigl(\zeta(2s)-1-2^{-2s}\bigr).
\]
Für $s\ge3/2$ gilt
\[
 q_s\le q_{3/2}
 =8\left(\zeta(3)-\frac98\right)
 <1.
\]
Die induzierte Rückkehrfolge ist in diesem Bereich also defekt: Die
Summe der gleichmäßigen Blockmajoranten ist strikt kleiner als $1$.
Ein diskretes defektes Erneuerungslemma für den Schwanz
$(m+1)^{-2s}$ würde daraus unmittelbar
$A_j(s)\ll_s j^{-2s}$ und damit die scharfe Ebenenasymptotik in diesem
Teilbereich liefern. Die noch zu schreibende Analyse besteht genau darin,
dieses Erneuerungslemma mit dem terminalen parabolischen Rest zu
formalisieren.
\end{remark}

\begin{proposition}[Defekte operatorielle Erneuerung]
Für jedes $s>1$ gilt
\[
 B_k(s)\ll_s(k+1)^{-2s}
 \qquad(k\ge1).
\]
Folglich erfüllt die einseitige Kernmasse
\[
 A_k(s)\ll_s(k+1)^{-2s}.
\]
\end{proposition}

\begin{proof}
Setze für $n\ge1$
\[
 \widehat R_n=R_{n-1},
 \qquad
 \widehat R=\sum_{n\ge1}\widehat R_n=\mathcal G_s.
\]
Nach den vorangehenden Propositionen gilt auf $C^1(Y)$
\[
 \widehat R_n=n^{-2s}\Pi_s+O_s(n^{-2s-1}),
 \qquad
 (I-\widehat R)^{-1}\ \,\text{existiert}.
\]
Wir verwenden die defekte operatorielle Erneuerungsabschätzung für
regulär variierende Operatorfolgen, wie sie etwa in der
\href{https://arxiv.org/abs/1008.4113}{Operator-Erneuerungstheorie von
Melbourne--Terhesiu} formuliert wird. Für
\[
 U_0=I,
 \qquad
 U_k=\sum_{\ell\ge1}
 \sum_{n_1+\cdots+n_\ell=k}
 \widehat R_{n_1}\cdots\widehat R_{n_\ell}
 \quad(k\ge1)
\]
liefert sie
\[
 \|U_k\|_{C^1(Y)\to C^1(Y)}
 \ll_s(k+1)^{-2s}.
\]

Die terminalen parabolischen Läufe werden durch die Funktionen
\[
 \tau_n(x)=(1+nx)^{-2s},
 \qquad n\ge0,
\]
erfasst. Sie erfüllen
$\|\tau_n\|_{C^1}\ll_s(n+1)^{-2s}$, und die eindeutige
Rückkehrzerlegung jeder Quotientenbahn liefert
\[
 B_k(s)
 =
 \left(
 \sum_{j=0}^kU_j\tau_{k-j}
 \right)(1).
\]
Da für $p=2s>1$
\[
 \sum_{j=0}^k
 (j+1)^{-p}(k-j+1)^{-p}
 \ll_p(k+1)^{-p},
\]
folgt die behauptete Schranke für $B_k$. Die Identität
$B_k=2A_k$ liefert die zweite Aussage.
\end{proof}

\begin{corollary}[Scharfe parabolische Ebenenasymptotik]
Für jedes $s>1$ gilt
\[
 \mathcal Z_k^{\mathrm{ref}}(s)
 \sim
 2^{s+1}\frac{\zeta(2s-1)}{\zeta(2s)}\,k^{-s}.
\]
\end{corollary}

\begin{proof}
Die vorangehende Proposition erfüllt das ausreichende Kernkriterium.
\end{proof}

\begin{remark}[Konkreter verbleibender Analyseschritt]
Die vorangehende Proposition kontrolliert die zwischenzeitliche zentrale
Dynamik bereits durch eine konvergente Resolvente. Die induzierten
Rückkehrblöcke tragen zugleich die erwartete parabolische Schwanzordnung
$m^{-2s}$. Offen bleibt daher nur noch ein operatorwertiges
Erneuerungslemma nach exakter Gesamtlänge: Es muss die Resolventenkontrolle
mit der asymptotischen Form
\[
 m^{2s}R_mf(x)\longrightarrow x^{-2s}f(1)
\]
der einzelnen Rückkehroperatoren verbinden. Dieser Schritt soll die
Schranke $A_j(s)\ll_s j^{-2s}$ im gesamten Bereich $s>1$ liefern.
\end{remark}

\subsection{Zweiparametrige Tiefenzeta}

Die Tiefenvariable kann unabhängig von der geometrischen Gewichtung
markiert werden. Für $\operatorname{Re}s>1$ und $|z|<1$ setzen wir
\[
 \mathscr D(s,z)
 =
 \sum_{k\ge0}z^k\mathcal Z_k^{\mathrm{ref}}(s)
 =
 \sum_{w\in\{L,R\}^*}
 z^{|w|}\left(\frac{2}{u_wv_w}\right)^s.
\]
Die Reihe konvergiert absolut und definiert eine in beiden Variablen
holomorphe Funktion. Bei $z=1$ erhält man die vollständige
Referenzzeta zurück:
\[
 \mathscr D(s,1)
 =
 2^s\frac{\zeta(s)^2}{\zeta(2s)}.
\]

Die Quotientendynamik liefert zudem eine Transferdarstellung. Mit
\[
 \psi_s(x)=x^{-s},
 \qquad
 (\mathcal L_sf)(x)
 =(1+x)^{-2s}
 \left(
 f\!\left(\frac{x}{1+x}\right)
 +f\!\left(\frac1{1+x}\right)
 \right)
\]
gilt, zunächst als absolut konvergente Neumannreihe,
\[
 \mathscr D(s,z)
 =
 2^s\sum_{k\ge0}z^k
 (\mathcal L_s^k\psi_s)(1).
\]
Damit ist die Tiefenmarkierung $z$ ein spektraler Parameter der
Stern--Brocot-Dynamik.

\begin{proposition}[Abelsche Tiefensingularität]
Für reelles $1<s<2$ gilt beim Grenzübergang $z\uparrow1$
\[
 \mathscr D(s,z)-\mathscr D(s,1)
 \sim
 C_s\Gamma(1-s)(1-z)^{s-1},
\]
wobei
\[
 C_s=
 2^{s+1}\frac{\zeta(2s-1)}{\zeta(2s)}.
\]
Für $s=2$ ersetzt der führende Term sich durch
\[
 \mathscr D(2,z)-\mathscr D(2,1)
 \sim
 C_2(1-z)\log(1-z).
\]
\end{proposition}

\begin{proof}
Die bewiesene Ebenenasymptotik
$\mathcal Z_k^{\mathrm{ref}}(s)\sim C_sk^{-s}$ und der
abelsche Satz für Potenzreihen regulär variierender Koeffizienten liefern
die Aussage. Der Modellterm ist $C_s\operatorname{Li}_s(z)$.
\end{proof}

\begin{remark}[Trennung von Tiefe und Arithmetik]
Die nichtanalytische Struktur in der Variablen $z$ bei $z=1$ wird von den
parabolischen Rückkehrzeiten erzeugt. Die arithmetische Struktur in der
Variablen $s$ ist dagegen bereits auf der Schnittebene $z=1$ durch
$\zeta(s)^2/\zeta(2s)$ sichtbar. Diese Trennung ist der Ausgangspunkt
für eine spätere Resonanzanalyse.
\end{remark}

\subsection{Parabolische Regularisierung und Resonanzprogramm}

Für komplexes $s$ und $|z|<1 betrachten wir die induzierte
operatorwertige Tiefenreihe
\[
 \mathcal G(s,z)
 =\sum_{m\ge0}z^{m+1}R_m(s).
\]
Die Rang-eins-Asymptotik der Rückkehroperatoren erlaubt eine kanonische
Abspaltung des parabolischen Teils. Für $m\ge1$ schreibe
\[
 R_m(s)=m^{-2s}\Pi_s+E_m(s).
\]
Dann gilt auf jedem kompakten Teilgebiet von $\operatorname{Re}s>0$
\[
 \|E_m(s)\|_{C^1\to C^1}
 \ll m^{-2\operatorname{Re}s-1}.
\]
Folglich besitzt
\[
 \mathcal E(s,z)
 =
 zR_0(s)+z\sum_{m\ge1}z^mE_m(s)
\]
eine bis $|z|=1$ stetige operatorwertige Fortsetzung für
$\operatorname{Re}s>0$, und
\begin{equation}\label{eq:polylog-regularization}
 \mathcal G(s,z)
 =
 z\operatorname{Li}_{2s}(z)\Pi_s+\mathcal E(s,z).
\end{equation}

\begin{proof}
Die Abschätzung für $E_m(s)$ ist die komplexe, auf kompakten Mengen
gleichmäßige Form der Rang-eins-Schwanzabschätzung. Da die Restreihe wie
$\sum_mm^{-2\operatorname{Re}s-1}$ majorisiert wird, konvergiert sie
in Operatornorm. Die Hauptreihe ist per Definition der Polylogarithmus.
\end{proof}

\begin{remark}[Drei getrennte Singularitätsmechanismen]
Formel \eqref{eq:polylog-regularization} trennt drei mögliche Quellen
nichtanalytischen Verhaltens:
\[
 \begin{array}{rcl}
 z=1
 &:&
 \operatorname{Li}_{2s}(z)
 \quad\text{(parabolische Rückkehrzeiten)},\\[2mm]
 I-\mathcal G(s,z)
 &:&
 \text{Nichtinvertierbarkeit der induzierten Resolvente
 (dynamische Resonanzen)},\\[2mm]
 z=1
 &:&
 \displaystyle
 \mathscr D(s,1)=2^s\frac{\zeta(s)^2}{\zeta(2s)}
 \quad\text{(arithmetische Singularitäten).}
 \end{array}
\]
Die dritte Zeile ist im Referenzmodell exakt. Insbesondere können
Nullstellen $\rho$ von $\zeta$ über den Nenner bei
$s=\rho/2$ zu Singularitäten der meromorphen Fortsetzung führen,
sofern keine Kürzung eintritt. Die erste und zweite Zeile stammen
dagegen aus der parabolischen Dynamik selbst.
\end{remark}

\paragraph{Nächster Resonanzsatz.}
Der nächste analytische Schritt besteht darin, einen geeigneten
Funktionsraum zu wählen, auf dem die Regularisierung
\eqref{eq:polylog-regularization} die Nichtinvertierbarkeit von
$I-\mathcal G(s,z)$ diskret erfasst. Erst dann kann die Lage dynamischer
Resonanzen mit den arithmetischen Singularitäten von
$\zeta(2s)^{-1}$ sinnvoll verglichen werden.

\subsection{Holomorpher Fredholmraum}

Setze
\[
 \Omega=D\!\left(1,\frac34\right)
\]
und bezeichne mit $\mathcal A(\Omega)$ den Banachraum der auf
$\Omega$ holomorphen und auf $\overline\Omega$ stetigen Funktionen,
versehen mit der Supremumsnorm.

\begin{proposition}[Holomorphe Fredholmfamilie]
Für jedes $m\ge0$ gilt
\[
 G_m(\overline\Omega)\Subset\Omega.
\]
Folglich ist $\mathcal G(s,z)$ für
$\operatorname{Re}s>1/2$ und $|z|\le1$ eine kompakte, in $s$
holomorphe Operatorfamilie auf $\mathcal A(\Omega)$. Für jedes feste
$z$ mit $|z|\le1$ besitzt
\[
 (I-\mathcal G(s,z))^{-1}
\]
eine meromorphe Fortsetzung in die Halbebene
$\operatorname{Re}s>1/2$. Ihre Nichtinvertibilitäten sind dort diskret.
Ferner ist die Halbebene
\[
 \operatorname{Re}s>1
\]
resonanzfrei.
\end{proposition}

\begin{proof}
Für $w\in\overline\Omega$ ist $\operatorname{Re}w\ge1/4$. Daher
gilt
\[
 |G_m(w)-1|
 =\frac{|w|}{|1+(m+1)w|}
 \le\frac{|w|}{|1+w|}
 <\frac34.
\]
Die letzte strikte Ungleichung folgt direkt aus
$|w-1|\le3/4$. Damit liegen die Bilder aller Zweige kompakt in
$\Omega$. Die Gewichte
$(1+(m+1)w)^{-2s}$ sind dort holomorph und auf kompakten
$s$-Mengen durch ein Vielfaches von
$(m+1)^{-2\operatorname{Re}s}$ beschränkt. Die Reihe für
$\mathcal G(s,z)$ konvergiert daher in Operatornorm für
$\operatorname{Re}s>1/2$ und ist als Normgrenzwert kompakter
gewichteter Kompositionsoperatoren kompakt.

Für reelles $\sigma>1$ liefert die zuvor bewiesene Resolvente
\[
 \|\mathcal G(\sigma,1)^n\|\ll
 \left(\frac49\right)^{n(\sigma-1)}.
\]
Für komplexes $s$ mit $\operatorname{Re}s=\sigma$ und $|z|\le1$
gilt auf dem reellen Intervall $Y$ punktweise
\[
 |\mathcal G(s,z)^nf|
 \le\mathcal G(\sigma,1)^n|f|.
\]
Liegt $f\in\mathcal A(\Omega)$ im Kern von
$I-\mathcal G(s,z)$, so folgt durch Iteration $f|_Y=0$, also wegen
der Holomorphie $f=0$. Da $I-\mathcal G(s,z)$ ein Fredholmoperator
vom Index $0$ ist, ist er in der Halbebene
$\operatorname{Re}s>1$ invertierbar. Der analytische Fredholm-Satz
liefert die meromorphe Fortsetzung und die Diskretheit der verbleibenden
Nichtinvertibilitäten.
\end{proof}

\begin{remark}[Dynamische Resonanzen]
Für festes $z$ nennen wir die diskreten Punkte in
$\operatorname{Re}s>1/2$, an denen
$I-\mathcal G(s,z)$ nicht invertierbar ist, die dynamischen
Resonanzen der induzierten Tangenzdynamik. Sie sind von den
arithmetischen Singularitäten der expliziten Schnittebene $z=1$
zunächst unabhängig definiert.
\end{remark}


\subsection{Resolventendarstellung der Tiefenzeta}

Die Fredholmfamilie erfasst die volle zweiparametrige Tiefenzeta, sobald
auch der letzte, nicht zur Rückkehr abgeschlossene parabolische Lauf
mitgeführt wird. Dazu setze
\[
 \eta(s,z;x)
 =
 \sum_{n\ge0}z^n x^{-s}(1+nx)^{-s}.
\]
Für $|z|<1$ konvergiert diese Reihe lokal gleichmäßig in $s$ als
$\mathcal A(\Omega)$-wertige Reihe. Sie ist somit eine ganze
$\mathcal A(\Omega)$-wertige Funktion von $s$.

\begin{proposition}[Exakte induzierte Resolventendarstellung]
Für $\operatorname{Re}s>1$ und $|z|\le1$ gilt
\begin{equation}\label{eq:induced-depth-resolvent}
 \mathscr D(s,z)
 =
 2^s\,
 \bigl[(I-\mathcal G(s,z))^{-1}\eta(s,z;\,\cdot)\bigr](1).
\end{equation}
Folglich besitzt $\mathscr D(s,z)$ für jedes feste $|z|<1$ eine
meromorphe Fortsetzung in die Halbebene
\[
 \operatorname{Re}s>\frac12.
\]
Ihre möglichen Pole liegen unter den dynamischen Resonanzen, das heißt
unter den Nichtinvertibilitäten von $I-\mathcal G(s,z)$.
\end{proposition}

\begin{proof}
Jede endliche Quotientenbahn, die bei $x=1\in Y$ beginnt, besitzt
eindeutig die Form
\[
 (T_0^{m_1}T_1)\cdots(T_0^{m_r}T_1)\,T_0^n,
 \qquad m_1,\ldots,m_r,n\ge0.
\]
Der $j$-te Rückkehrblock trägt genau die Länge $m_j+1$ und wird durch
$R_{m_j}(s)$ erfasst. Für den abschließenden Lauf gilt
\[
 T_0^n(x)=\frac{x}{1+nx},
 \qquad
 \prod_{\ell=0}^{n-1}
  \bigl(1+T_0^\ell(x)\bigr)^{-2s}
 \,\psi_s(T_0^n x)
 =
 x^{-s}(1+nx)^{-s}.
\]
Die Markierung der Gesamtlänge liefert daher nach Summation über die
Rückkehrblöcke und den terminalen Lauf
\[
 \frac{\mathscr D(s,z)}{2^s}
 =
 \sum_{r\ge0}\mathcal G(s,z)^r\eta(s,z;\,\cdot)(1).
\]
Im anfänglichen Bereich $\operatorname{Re}s>1$ ist die ursprüngliche
Wortreihe auch für $|z|=1$ absolut konvergent. Die Zerlegung darf daher nach Zahl und
Längen der Rückkehrblöcke umgeordnet werden. Ihre Summe löst
\[
 (I-\mathcal G(s,z))F=\eta(s,z;\,\cdot).
\]
Wegen der dortigen Invertierbarkeit von $I-\mathcal G(s,z)$ ist
$F=(I-\mathcal G(s,z))^{-1}\eta(s,z;\,\cdot)$, und
\eqref{eq:induced-depth-resolvent} folgt. Die Fredholmfortsetzung der
Resolvente aus der vorangehenden Proposition und die Ganzeigenschaft
von $\eta$ liefern die meromorphe Fortsetzung. Pole können nur dort
entstehen, wo die Resolvente einen Pol besitzt; sie können im Einzelfall
durch die terminale Beobachtung weggekürzt werden.
\end{proof}


\subsection{Terminale parabolische Regularisierung}

Auch der terminale Lauf besitzt eine kanonische Polylogarithmus-
Abspaltung. Für $n\ge1$ schreiben wir
\[
 x^{-s}(1+nx)^{-s}
 =
 x^{-2s}n^{-s}
 +
 x^{-2s}n^{-s}
 \left[
 \left(1+\frac1{nx}\right)^{-s}-1
 \right].
\]
Daher gilt die exakte Zerlegung
\begin{equation}\label{eq:terminal-polylog-regularization}
 \eta(s,z;x)
 =
 x^{-2s}\operatorname{Li}_s(z)+\mathcal H(s,z;x),
\end{equation}
wobei
\[
 \mathcal H(s,z;x)
 =
 x^{-s}
 +
 x^{-2s}\sum_{n\ge1}z^n n^{-s}
 \left[
 \left(1+\frac1{nx}\right)^{-s}-1
 \right].
\]

\begin{proposition}[Regulärer terminaler Rest]
Die Funktion $\mathcal H(s,z;\,\cdot)$ ist für
$\operatorname{Re}s>0$ holomorph in $s$, holomorph in $z$ für
$|z|<1$ und setzt sich als $\mathcal A(\Omega)$-wertige Funktion
stetig nach $|z|\le1$ fort. Für reelles $1<s<2$ gilt in der
$\mathcal A(\Omega)$-Norm
\[
 \eta(s,z;\,\cdot)-\eta(s,1;\,\cdot)
 \sim
 \Gamma(1-s)(1-z)^{s-1}\,(\,\cdot\,)^{-2s}
 \qquad(z\uparrow1).
\]
Bei $s=2$ lautet der führende nichtanalytische Term
\[
 (1-z)\log(1-z)\,(\,\cdot\,)^{-4}.
\]
\end{proposition}

\begin{proof}
Auf $\overline\Omega$ ist $x$ gleichmäßig von $0$ getrennt. Für $s$
in einer kompakten Teilmenge von $\operatorname{Re}s>0$ gilt daher
gleichmäßig in $x$
\[
 n^{-s}
 \left[
 \left(1+\frac1{nx}\right)^{-s}-1
 \right]
 =
 O\!\left(n^{-\operatorname{Re}s-1}\right).
\]
Die Reihe für $\mathcal H$ konvergiert folglich bis $|z|=1$ in der
Norm von $\mathcal A(\Omega)$. Die behaupteten Regularitätseigenschaften
folgen aus der lokal gleichmäßigen Konvergenz. Für $s>1$ ist der Rest
sogar einmal stetig nach $z$ differenzierbar am Punkt $z=1$. Die
nichtanalytischen Leitterme stammen deshalb allein von
$\operatorname{Li}_s(z)$ in \eqref{eq:terminal-polylog-regularization};
die bekannten Polylogarithmus-Asymptotiken liefern die Aussagen.
\end{proof}

\begin{corollary}[Der terminale Pol bei $s=1$]
Auf der Randebene $z=1$ gilt zunächst für $\operatorname{Re}s>1$,
und danach meromorph in $s$,
\[
 \eta(s,1;x)
 =
 x^{-2s}\zeta\!\left(s,\frac1x\right),
\]
wobei rechts die Hurwitz-Zeta-Funktion steht. Insbesondere besitzt
der terminale Anteil bei $s=1$ einen einfachen Pol mit Residuum
\[
 \operatorname*{Res}_{s=1}\eta(s,1;x)=x^{-2}.
\]
\end{corollary}

\begin{proof}
Aus $1+nx=x(n+x^{-1})$ folgt
\[
 \eta(s,1;x)
 =
 x^{-2s}\sum_{n\ge0}\left(n+\frac1x\right)^{-s}.
\]
Die Aussage folgt aus der Standardfortsetzung der Hurwitz-Zeta-Funktion.
\end{proof}

\begin{proposition}[Explizite Randresonanz]
Am Punkt $(s,z)=(1,1)$ gilt
\[
 \mathcal G(1,1)\,h=h,
 \qquad h(x)=x^{-1}.
\]
Insbesondere ist $s=1$ auf der Randebene $z=1$ tatsächlich eine
dynamische Resonanz.
\end{proposition}

\begin{proof}
Da $|G_m'(x)|=(1+(m+1)x)^{-2}$, erhält man
\[
 \begin{aligned}
 (\mathcal G(1,1)h)(x)
 &=
 \sum_{m\ge0}
 \frac{1}{(1+(m+1)x)^2}\,
 \frac{1+(m+1)x}{1+mx}\\
 &=
 \sum_{m\ge0}
 \frac1{(1+mx)(1+(m+1)x)}\\
 &=
 \frac1x\sum_{m\ge0}
 \left(
 \frac1{1+mx}-\frac1{1+(m+1)x}
 \right)
 =\frac1x.
 \end{aligned}
\]
\end{proof}


\subsection{Resolventennachweis des Doppelpols}

Die induzierte Dynamik ist nach einer elementaren Koordinatenänderung
der klassische Gauß-Transferoperator. Setze dazu
\[
 K(u)=\frac1{1+u},
 \qquad
 S_n(u)=\frac1{n+u}\quad(n\ge1),
\]
und definiere
\[
 (U_sf)(u)=(1+u)^{-2s}f(K(u)).
\]
Dann gilt die exakte Konjugationsformel
\[
 U_s\,\mathcal G(s,1)\,U_s^{-1}=\mathcal L_s^{\mathrm G},
 \qquad
 (\mathcal L_s^{\mathrm G}F)(u)
 =
 \sum_{n\ge1}(n+u)^{-2s}F(S_n(u)).
\]
Die Rückkehrdynamik ist also keine neue spektrale Klasse: Sie ist die
Gauß-Dynamik in der Quotientenkoordinate $u=x^{-1}-1$.

\begin{proposition}[Spektraler Gauß-Input]
In einer Umgebung von $s=1$ besitzt
$\mathcal L_s^{\mathrm G}$ auf einem geeigneten holomorphen
Transferraum einen eindeutig bestimmten, holomorph von $s$ abhängigen
Eigenwert $\lambda(s)$ mit
\[
 \lambda(1)=1.
\]
Dieser Eigenwert ist bei $s=1$ algebraisch einfach; der verbleibende
Spektralteil ist von $1$ getrennt. Dieselbe Aussage gilt daher für
$\mathcal G(s,1)$.
\end{proposition}

\begin{proof}
Dies ist der klassische Spektralsatz für den Gauß--Kuzmin--Wirsing-
Operator auf holomorphen Funktionsräumen. Die hier verwendete
Operatorfamilie ist exakt die Familie mit Gewichten
$(n+u)^{-2s}$. Eine passende holomorphe Formulierung findet sich etwa
bei \href{https://maths.qmul.ac.uk/~omj/fm6.pdf}
{Jenkinson--González--Urbański}; die Konjugationsformel überträgt sie
auf unsere induzierte Familie.
\end{proof}

\begin{proposition}[Ableitung des führenden Eigenwerts]
Für den führenden Eigenwert gilt
\[
 \lambda'(1)=-\frac{\pi^2}{6\log2}.
\]
\end{proposition}

\begin{proof}
Für $\mathcal L_1^{\mathrm G}$ sind rechter und linker Eigenvektor
\[
 \rho(u)=\frac1{1+u},
 \qquad
 \ell(F)=\int_0^1F(u)\,du,
 \qquad
 \ell(\rho)=\log2.
\]
Analytische Störungstheorie liefert daher
\[
 \lambda'(1)
 =
 \frac{\ell\bigl((\mathcal L_1^{\mathrm G})'\rho\bigr)}
      {\ell(\rho)}.
\]
Nach der Variablentransformation $v=(n+u)^{-1}$ in jedem Zweig folgt
\[
 \ell\bigl((\mathcal L_1^{\mathrm G})'\rho\bigr)
 =
 2\int_0^1\frac{\log v}{1+v}\,dv
 =
 -\frac{\pi^2}{6}.
\]
In der letzten Gleichheit wurde
$(1+v)^{-1}=\sum_{j\ge0}(-v)^j$ integriert. Division durch $\log2$
ergibt die Behauptung.
\end{proof}

\begin{theorem}[Resolvententheoretische Herleitung des Doppelpols]
Die Referenz-Tiefenzeta besitzt bei $s=1$ die Entwicklung
\[
 \mathscr D(s,1)
 =
 \frac{12}{\pi^2}(s-1)^{-2}
 +O\!\left((s-1)^{-1}\right).
\]
Insbesondere wird ihr Doppelpol bei $s=1$ allein aus terminalem
parabolischem Pol und einfacher dynamischer Resonanz hergeleitet.
\end{theorem}

\begin{proof}
Sei $h(x)=x^{-1}$ und $\nu(f)=\int_{1/2}^1f(x)\,dx$. Da die
Rückkehrintervalle $G_m(Y)$ $Y$ bis auf Endpunkte zerlegen, ist
$\nu$ linker Eigenvektor von $\mathcal G(1,1)$ und
\[
 \nu(h)=\log2.
\]
Die spektrale Projektion zum Eigenwert $1$ ist folglich
\[
 Pf=h\,\frac{\nu(f)}{\log2}.
\]
Aus der Einfachheit und der berechneten Ableitung folgt für die
Resolvente
\[
 (I-\mathcal G(s,1))^{-1}
 =
 -\frac{P}{\lambda'(1)}\frac1{s-1}+O(1).
\]
Andererseits besitzt der terminale Anteil nach dem vorangehenden
Korollar die Form
\[
 \eta(s,1;\,\cdot)
 =
 \frac{q}{s-1}+O(1),
 \qquad q(x)=x^{-2}.
\]
Da
\[
 (Pq)(1)
 =
 \frac{\int_{1/2}^1x^{-2}\,dx}{\log2}
 =
 \frac1{\log2},
\]
liefert \eqref{eq:induced-depth-resolvent}
\[
 \mathscr D(s,1)
 =
 -\frac{2(Pq)(1)}{\lambda'(1)}(s-1)^{-2}
 +O\!\left((s-1)^{-1}\right)
 =
 \frac{12}{\pi^2}(s-1)^{-2}
 +O\!\left((s-1)^{-1}\right).
\]
\end{proof}

\begin{remark}[Kollision am Rand]
Im offenen Bereich $|z|<1$ sind terminale und dynamische
Singularitäten getrennt. Auf $z=1$ treffen sie bei $s=1$ zusammen:
Der terminale Anteil hat dort einen einfachen Hurwitz-Zeta-Pol, und
zugleich besitzt die induzierte Dynamik die Eigenfunktion $x^{-1}$ zum
Eigenwert $1$. Der vorangehende Satz zeigt, dass diese beiden einfachen
Mechanismen mit nichtverschwindender Kopplung genau den Doppelpol
\[
 \mathscr D(s,1)\sim\frac{12}{\pi^2}(s-1)^{-2}
\]
erzeugen.
\end{remark}

\begin{remark}[Präzise Trennung der Mechanismen]
Für $|z|<1$ ist die terminale Funktion $\eta(s,z;\,\cdot)$ ganz in
$s$. In diesem offenen Tiefenbereich entstehen Singularitäten der
Fortsetzung daher ausschließlich dynamisch. Die arithmetische
Singularitätsstruktur
\[
 \mathscr D(s,1)=2^s\frac{\zeta(s)^2}{\zeta(2s)}
\]
tritt erst an der Randebene $z=1$ hinzu. Dies ist eine strikte
Trennung, keine bloße Analogie.
\end{remark}

\section{Offene Fragen}

\begin{enumerate}[leftmargin=*,label=\arabic*.]
\item Lässt sich der effektive lokale Kreis-Zählsatz auf dem kompakten
Formraum unserer Wurzellücken mit einem einheitlichen Exponenten und einer
einheitlichen Restkonstante formulieren? Dann verfeinert sich
\eqref{eq:global-main-asymptotic} zu \eqref{eq:global-exact-count}, und
die Abel-Grenze wird zu einem echten meromorphen Residuum.
 \item Gilt für jede Wurzellücke die scharfe Ebenenasymptotik
 \[
  \mathcal Z_{n,k}(s)\sim C_{n,s}k^{-s}
 \]
 für $s>1$?
\item Welche Hausdorff-Dimensionen besitzen die Niveauklassen
\[
 \left[\frac{\log2}{2\log\varphi},+\infty\right]
\]
und wie verändern sie sich zwischen dem linken Rand, dem typischen Wert
und der parabolischen Kusp? Die Realisierung aller endlichen Werte durch
nichtperiodische Adressen ist bereits gesichert.
 \item Besitzt der kanonische Quotient $X/{\sim_\pi}$ eine endliche oder
 niedrigdimensionale geometrische Realisierung?
 \item Gilt die Kuspvermutung $\Haus E(+\infty)=1/2$? Die Untergrenze ist
 nun bewiesen; offen ist allein die passende globale Oberschranke.
\end{enumerate}

\section{Ausblick}

Die Ausgangsfrage hat zwei gleichwertige, aber metrisch verschiedene
Antworten. Globale Halbierung erzeugt ein lineares dyadisches Raster.
Tangenz erzeugt eine Möbius-Dynamik, in der gleiche symbolische Information
pro Schritt geometrisch sehr ungleich skaliert wird. Genau diese Verzerrung
eröffnet die Zeta- und Multifraktalfragen. Für die kanonische Randfamilie
an der Halbierungsgeraden ist das Wörter--Lücken-Wörterbuch nun vollständig
geometrisch begründet; Exponent und Zählordnung sind dort Theoreme. Auch die
Außenkreisfamilie ist durch eine Cayley-Transformation auf denselben
Mechanismus zurückgeführt. Für die vollständige Packung liegt nun eine
eindeutige Zählkonvention vor; kritischer Exponent, Zählordnung und
Zählhauptkonstante sind bestimmt. Der nächste sinnvolle Schritt ist nun die
uniform effektive lokale Zählung, die einen globalen Potenzrestterm und eine
meromorphe Zeta-Fortsetzung liefern würde. Für das Multifraktalspektrum ist
der linke Rand nun scharf bestimmt, und jeder endliche Wert ist auch
nichtperiodisch realisierbar; periodische Werte sind darüber hinaus dicht.
Für die parabolische Kusp ist die Untergrenze $1/2$ bewiesen. Offen
bleiben die Ebenensummen, die Hausdorff-Dimensionen der endlichen
Niveauklassen und die globale Oberschranke für die Kusp.

\end{document}
